% ============ 第一章 I 局部类域论: Lubin-Tate 理论 ============
% 原书 Chapter I, pp.19--52 (局部类域论: Lubin-Tate 理论)
% 编号: 定理/推论/引理/命题/定义/例/注记/习题 共享计数器, 按节编号(1.1, 1.2, ...)
% 节内说明性编号 1.6--1.8, 1.11, 1.12, 1.14, 1.15 亦占该计数器;
% 其中 1.9, 1.10 为子节标题, 不占计数器, 故在 1.11 前将计数器置为 10.
\chapter{局部类域论: Lubin-Tate 理论}\label{ch:I}%
\index{jubuleiyulun@局部类域论(local class field theory)}\index{lubintatexingshiqunlü@Lubin-Tate 形式群律(Lubin-Tate formal group)}

局部类域论对局部域的 Abel 扩张进行分类. 换一个角度, 它描述了整体 Artin 映射的局部分量.

所谓局部域, 我是指对一个非平凡绝对值而言局部紧的域 $K$. 于是 (ANT, 7.49) 它或者是

(a) 对某个 $p$, $\Q_p$ 的有限扩张;

(b) 具有 $p$ 个元素的域上的 Laurent 级数域 $\F_p((T))$ 的有限扩张; 或者

(c) $\R$ 或 $\CC$ (阿基米德情形).

当 $K$ 非阿基米德 时, $K$ 中的整数环 (又称赋值环) 记为 $\mathcal{O}_K$ (或 $A$), 其极大理想记为 $\mathfrak{m}_K$ (或简记 $\mathfrak{m}$), 其单位群记为 $\mathcal{O}_K^{\times}$ 或 $U_K$. $\mathfrak{m}$ 的一个生成元称为 $K$ 的素元 (或称一致素元, 或局部一致化参数). 若 $\pi$ 是 $K$ 的素元, 则 $K^{\times}$ 中每个元素都可唯一写成 $a=u\pi^m$ 的形式, 其中 $u\in\mathcal{O}_K^{\times}$, $m\in\Z$. 我们定义
\[
\mathrm{ord}_K(a)=m.
\]
存在一个正合列
\[
0\to U_K\to K^{\times}\xrightarrow{\mathrm{ord}_K}\Z\to 0.
\]
$K$ 的剩余类域 $k$ 有 $q$ 个元素, 其特征为 $p$. $K$ 上的标准绝对值定义为 $|a|=q^{-\mathrm{ord}_K(a)}$.

我们以 $K^{\mathrm{al}}$ 记 $K$ 的一个固定的代数闭包 (当 $K$ 的特征 $p\neq 0$ 时, 指可分代数闭包), “$K$ 的扩张”意即 “含 $K$ 的 $K^{\mathrm{al}}$ 的子域”. $\mathrm{ord}_K$ 与 $|\cdot|$ 都可唯一扩张到 $K^{\mathrm{al}}$ ($\mathrm{ord}_K$ 的扩张取值于 $\Q$).

\section{主定理的陈述}\label{sec:I.1}

两个 $K$ 的有限 Abel 扩张的复合仍是 $K$ 的有限 Abel 扩张 (FT, 3.20). 因此所有有限 Abel 扩张在 $K^{\mathrm{al}}$ 中的并 $K^{\mathrm{ab}}$ 是一个无限 Abel 扩张, 其 Galois 群是 $\Gal(K^{\mathrm{al}}/K)$ 模去其换位子群之闭包所得的商群. (见 FT 第 7 章及本章附录中关于无限扩张的 Galois 理论.)

设 $L$ 是 $K$ 的有限非分歧扩张. 则 $L/K$ 是 Galois 的, 且 $\Gal(L/K)$ 在 $\mathcal{O}_L$ 上的作用定义了一个同构 $\Gal(L/K)\simeq\Gal(l/k)$, 其中 $l$ 是 $L$ 的剩余类域. 因此 $\Gal(L/K)$ 是循环群, 由唯一的元素 $\sigma$ 生成, 该元素对所有 $\alpha\in\mathcal{O}_L$ 满足 $\sigma\alpha\equiv\alpha^q\bmod\mathfrak{m}_L$. 这个 $\sigma$ 称为 $\Gal(L/K)$ 的 \emph{Frobenius 元素}, 记为 $\Frob_{L/K}$. 见 ANT, 7.54.

\begin{theorem}[局部互反律]\label{thm:1.1}
对每个非阿基米德局部域 $K$, 存在唯一的同态
\[
\phi_K\colon K^{\times}\to\Gal(K^{\mathrm{ab}}/K)
\]
满足以下性质:

(a) 对 $K$ 的每个素元 $\pi$ 以及 $K$ 的每个有限非分歧扩张 $L$, $\phi_K(\pi)$ 在 $L$ 上的作用等于 $\Frob_{L/K}$;

(b) 对 $K$ 的每个有限 Abel 扩张 $L$, $\Nm_{L/K}(L^{\times})$ 包含在 $a\mapsto\phi_K(a)|_L$ 的核中, 且 $\phi_K$ 诱导出一个同构
\[
\phi_{L/K}\colon K^{\times}/\Nm_{L/K}(L^{\times})\to\Gal(L/K).
\]
特别地,
\[
(K^{\times}:\Nm_{L/K}(L^{\times}))=[L:K].
\]
\end{theorem}

记 $\Nm_{L/K}(L^{\times})$ 为 $\Nm(L^{\times})$. 陈述 (b) 表明, 对每个有限 Abel 扩张 $L/K$, 映射 $\phi_K$ 可分解为
\[
\begin{tikzcd}
K^{\times}\ar[r,"\phi_K"]\ar[d] & \Gal(K^{\mathrm{ab}}/K)\ar[d,"\tau\mapsto\tau|_L"]\\
K^{\times}/\Nm(L^{\times})\ar[r,"\simeq","\phi_{L/K}"'] & \Gal(L/K).
\end{tikzcd}
\]
我们称 $\phi_K$ 与 $\phi_{L/K}$ 为 $K$ 与 $L/K$ 的\emph{局部 Artin 映射}. 它们也常被称为\emph{局部互反律映射}, 记为 $\mathrm{rec}_K$ 与 $\mathrm{rec}_{L/K}$; 而 $\phi_{L/K}$ 也常称为\emph{范数剩余映射或范数剩余符号}, 记作 $a\mapsto(a,L/K)$.

形如 $\Nm(L^{\times})$ ($L$ 跑遍 $K$ 的有限 Abel 扩张) 的 $K^{\times}$ 的子群称为 $K^{\times}$ 中的\emph{范数群}\index{fanshuqun@范数群(norm group)}.

\begin{corollary}\label{cor:1.2}
设 $K$ 为非阿基米德局部域, 并假设存在同态 $\phi\colon K^{\times}\to\Gal(K^{\mathrm{ab}}/K)$ 满足定理的条件 (a) 与 (b).

(a) 映射 $L\mapsto\Nm(L^{\times})$ 是从 $K$ 的有限 Abel 扩张的集合到 $K^{\times}$ 中范数群的集合上的双射.

(b) $L\subset L'\iff \Nm(L^{\times})\supset\Nm(L'^{\times})$.

(c) $\Nm((L\cdot L')^{\times})=\Nm(L^{\times})\cap\Nm(L'^{\times})$.

(d) $\Nm((L\cap L')^{\times})=\Nm(L^{\times})\cdot\Nm(L'^{\times})$.\erratum{原书此处记作 $\Nm((L\cap L)^{\times})$, 依 Milne errata (v4.03) 应理解为 $L\cap L'$.}

(e) 每个包含某个范数群的 $K^{\times}$ 的子群自身也是范数群.
\end{corollary}

\begin{proof}
注意范次的传递性
\[
\Nm_{L'/K}=\Nm_{L/K}\circ\Nm_{L'/L}
\]
表明 $L\subset L'\implies \Nm(L^{\times})\supset\Nm(L'^{\times})$. 因此
\[
\Nm((L\cdot L')^{\times})\subset\Nm(L^{\times})\cap\Nm(L'^{\times}).
\]
反之, 若 $a\in\Nm(L^{\times})\cap\Nm(L'^{\times})$, 则
\[
\phi_{L/K}(a)=1=\phi_{L'/K}(a).
\]
但 $\phi_{LL'/K}(a)|_L=\phi_{L/K}(a)$, 且 $\phi_{LL'/K}(a)|_{L'}=\phi_{L'/K}(a)$. 由于映射
\[
\sigma\mapsto(\sigma|_L,\sigma|_{L'})\colon\Gal(LL'/K)\to\Gal(L/K)\times\Gal(L'/K)
\]
是单的 (FT, 3.20), 这就推出 $\phi_{LL'/K}(a)=1$, 从而 $a\in\Nm((L\cdot L')^{\times})$. 这证明了 (c); 现在来完成 (b) 的证明. 当 $\Nm(L^{\times})\supset\Nm(L'^{\times})$ 时, (c) 成为
\[
\Nm((LL')^{\times})=\Nm(L'^{\times}).
\]
由于一个范数群的指数等于定义它的 Abel 扩张的次数 (\ref{thm:1.1}b), 且 $LL'\supset L'$, 这表明 $LL'=L'$, 从而 $L'\supset L$.

现在证明 (a). 按定义, 映射 $L\mapsto\Nm(L^{\times})$ 是满的, 而由 (b) 知它是单的.

接着证明 (e). 设 $N$ 是一个范数群, 比如 $N=\Nm(L^{\times})$, 并设 $I\supset N$. 令 $M$ 为 $\phi_{L/K}(I)$ 的固定域, 于是 $\phi_{L/K}$ 把 $I/N$ 同构地映到 $\Gal(L/M)$ 上. 考虑交换图表
\[
\begin{tikzcd}
K^{\times}\ar[r,"\phi_{L/K}"]\ar[d,equal] & \Gal(L/K)\ar[d]\\
K^{\times}\ar[r,"\phi_{M/K}"'] & \Gal(M/K)
\end{tikzcd}
\]
其中右侧竖向箭头为限制映射（左侧为恒等映射）. $\phi_{M/K}$ 的核是 $\Nm(M^{\times})$. 另一方面, 映射
\[
K^{\times}\to\Gal(L/K)\to\Gal(M/K)
\]
的核是 $\phi_{L/K}^{-1}(\Gal(L/M))$, 按 Galois 理论它等于 $I$.

最后证明 (d). 在 (a) 中的两个集合之间有一个反序的双射. 由于 $L\cap L'$ 是同时含于 $L$ 与 $L'$ 中的 $K$ 的最大扩张, 而 $\Nm(L^{\times})\cdot\Nm(L'^{\times})$ 是包含 $\Nm(L^{\times})$ 与 $\Nm(L'^{\times})$ 的最小子群 (并且由 (e) 知它是范数群), 两者必相对应.
\end{proof}

为了对 $K$ 的 Abel 扩张进行分类, 还需确定范数群. 下一个引理表明 (在假定定理\ref{thm:1.1} 的前提下) 它们都是开集.

\begin{lemma}\label{lem:1.3}
设 $L$ 是 $K$ 的扩张. 若 $\Nm(L^{\times})$ 在 $K^{\times}$ 中指数为有限, 则它是开集.
\end{lemma}

\begin{proof}
$L$ 中的单位群 $U_L$ 是紧的, 故 $\Nm(U_L)$ 在 $K^{\times}$ 中闭. 考察赋值可知
\[
\Nm(L^{\times})\cap U_K=\Nm(U_L),
\]
于是 $U_K/\Nm(U_L)\hookrightarrow K^{\times}/\Nm(L^{\times})$. 因此 $\Nm(U_L)$ 在 $U_K$ 中指数为有限且闭, 从而在 $U_K$ 中开; 而 $U_K$ 自身在 $K^{\times}$ 中开. 于是 $\Nm(L^{\times})$ 含有 $K^{\times}$ 的一个开子群, 故自身为开.\footnote{这里用到了: 拓扑群中有限指数的闭子群是开的 (因为其补集是有限多个陪集的并, 而陪集也是闭的); 以及: 一个子群若含有一个开子群, 则它自身为开 (因为它是该开子群若干陪集的并)。}
\end{proof}

\begin{theorem}[局部存在定理]\label{thm:1.4}
$K^{\times}$ 中的范数群恰为有限指数的开子群.
\end{theorem}

\begin{corollary}\label{cor:1.5}
推论\ref{cor:1.2} 成立, 其中 “范数群” 替换为 “有限指数的开子群”.
\end{corollary}

在概述定理\ref{thm:1.1} 与 \ref{thm:1.4} 的证明之前, 我先给出关于这两个定理的一些补充.

\begin{npara}\label{para:1.6}
推论\ref{cor:1.5} 对阿基米德局部域也成立. $\R$ 的 Abel 扩张是 $\R$ 与 $\CC$, 它们的范数子群分别是 $\R^{\times}$ 与 $\R_{>0}$. 设 $H$ 是 $\R^{\times}$ 中有限指数的子群. 则对某个 $m$ 有 $H\supset\R^{\times m}$, 而依据介值定理, $\R^{\times m}=\R^{\times}$ 或 $\R_{>0}$ 分别视 $m$ 为奇或偶而定. 因此 $\R^{\times}$ 与 $\R_{>0}$ 是 $\R^{\times}$ 中仅有的两个有限指数子群. 唯一的同构
\[
\R^{\times}/\R_{>0}\to\Gal(\CC/\R)
\]
称为 $\R$ 的局部 Artin 映射.
\end{npara}

\begin{npara}\label{para:1.7}
当 $K$ 的特征为零时, $K^{\times}$ 的每个有限指数子群都是开的. 要证明这一点, 注意一个有限指数的子群 $H$ 将包含某个 $m$ 的 $K^{\times m}$; 而对满足 $|1-a|<|m|^2$ 的每个 $a\in\mathcal{O}_K^{\times}$, Newton 引理 (ANT, 7.32) 用于 $X^m-a$ 表明 $a$ 可写成 $a=u^m$, 其中 $u\in1+\mathfrak{m}$. 于是 $H$ 包含 $K^{\times}$ 中 $1$ 的一个开邻域, 而由于它是群, 这就推出它是开的.

当 $K$ 的特征 $p\neq 0$ 时, 并非 $K^{\times}$ 的每个有限指数子群都是开的. \cite{Weil1967}, II 3, 命题 10 表明 $1+\mathfrak{m}\approx\prod_{\N}\Z_p$ (以 $\N$ 为指标的 $\Z_p$ 的拷贝之积), 由此推出 $K^{\times}$ 有一个同构于 $\prod_{\N}\F_p$ 的商群. 子群 $\bigoplus_{\N}\F_p$ 在 $\prod_{\N}\F_p$ 中稠密. 因此 $\prod_{\N}\F_p$ 中每个包含 $\bigoplus_{\N}\F_p$ 的真子群都不闭. 选取这样一个有限指数的子群, 令 $U$ 为它在 $1+\mathfrak{m}$ 中的原像. 则 $U\cdot\pi^{\Z}$ 在 $K^{\times}$ 中指数为有限但不闭 (从而不开). 参见 FT 7.26.
\end{npara}

\begin{npara}\label{para:1.8}
$K$ 的两个有限非分歧扩张的复合仍是非分歧的, 因此所有有限非分歧扩张在 $K^{\mathrm{al}}$ 中的并 $K^{\mathrm{un}}$ 是 $K$ 的一个非分歧扩张. $K^{\mathrm{un}}$ 的剩余类域 $\bar{k}$ 是 $K$ 的剩余类域 $k$ 的代数闭包.

$K^{\mathrm{un}}$ 的每个固定 $K$ 的自同构 $\sigma$ 保持 $K^{\mathrm{un}}$ 上的绝对值 $|\cdot|$, 从而诱导出 $\bar{k}/k$ 的一个自同构 $\bar{\sigma}$. 映射 $\sigma\mapsto\bar{\sigma}$ 是一个同构
\[
\Gal(K^{\mathrm{un}}/K)\to\Gal(\bar{k}/k).
\]
因此存在唯一的元素 $\Frob_K\in\Gal(K^{\mathrm{un}}/K)$ 在 $\bar{k}$ 上诱导映射 $x\mapsto x^q$, 且映射 $\alpha\mapsto\Frob_K^{\alpha}\colon\hat{\Z}\to\Gal(K^{\mathrm{un}}/K)$ 是拓扑群的同构. 定理\ref{thm:1.1} 的条件 (a) 可以重新表述为: 对 $K$ 的每个素元 $\pi$, $\phi_K(\pi)$ 在 $K^{\mathrm{un}}$ 上的作用等于 $\Frob_K$. 特别地, 对任何单位 $u\in\mathcal{O}_K^{\times}$, $\phi_K(\pi)$ 与 $\phi_K(\pi u)$ 在 $K^{\mathrm{un}}$ 上有相同的作用, 因此 $\phi_K(u)$ 在 $K^{\mathrm{un}}$ 上起恒等映射的作用.
\end{npara}

\subsection*{1.9 群}

\begin{equation*}
U_K\supset 1+\mathfrak{m}\supset\cdots\supset 1+\mathfrak{m}^n\supset\cdots
\tag{7}\label{eqn:i.1.9.filt}
\end{equation*}
构成 $U_K$ 中 $1$ 处的一组基本邻域系. 设 $L$ 是 $K$ 的有限 Abel 扩张. 若 $L/K$ 非分歧, 则由 (\ref{para:1.8}) 知 $\Ker(\phi_{L/K})\supset U_K$, 我们说 $L/K$ 的\emph{导子}\index{daozi@导子(conductor)}为 $0$. 否则, 使得 $\Ker(\phi_{L/K})$ 包含 $1+\mathfrak{m}^f$ 的最小 $f$ 称为 $L/K$ 的\emph{导子}. 当且仅当导子为 $0$ 时 $L/K$ 非分歧, 当且仅当导子 $\le 1$ 时 $L/K$ 为顺分歧 (tamely ramified). [注意对 $i\ge 1$, 商 $1+\mathfrak{m}^i/1+\mathfrak{m}^{i+1}$ 是 $p$-群; 因此若 $L/K$ 顺分歧, 则 $\phi_{L/K}|_{U_K}$ 经过 $U_K/(1+\mathfrak{m})$. 更精细的论述见 \ref{para:1.11}.]

\subsection*{1.10 同态}

\[
\phi_{L/K}\colon K^{\times}/\Nm(L^{\times})\to\Gal(L/K)
\]
当 $L$ 跑遍 $K$ 的有限 Abel 扩张 (按包含排序) 时构成一个逆向系. 取极限, 我们得到一个同构
\[
\hat{\phi}_K\colon\widehat{K^{\times}}\to\Gal(K^{\mathrm{ab}}/K),
\]
其中 $\widehat{K^{\times}}$ 是 $K^{\times}$ 关于如下拓扑的完备化: 以范数群作为 $1$ 处的基本邻域系. $K^{\times}$ 上的这个拓扑称为\emph{范数拓扑}\index{fanshutopu@范数拓扑(norm topology)}. 按定理\ref{thm:1.4}, 范数群就是 $K^{\times}$ 中有限指数的开子群. 选取一个素元 $\pi$, 并写
\begin{equation*}
K^{\times}=U_K\cdot\pi^{\Z}\simeq U_K\times\Z.
\tag{8}
\end{equation*}
子群
\[
(1+\mathfrak{m}^n)\cdot\langle\pi^m\rangle\simeq(1+\mathfrak{m}^n)\times m\Z
\]
构成 $1$ 处的一组基本邻域系. 特别地, 对范数拓扑而言 $U_K$ 在 $K^{\times}$ 中不开, 因此该拓扑比 $K^{\times}$ 上通常的拓扑更粗. 完备化后, 上式成为
\[
\widehat{K^{\times}}=U_K\cdot\pi^{\hat{\Z}}\simeq U_K\times\hat{\Z},
\]
其中 $\hat{\Z}$ 是 $\Z$ 关于有限指数子群所定义拓扑的完备化. 更规范地说, 恒等映射
\[
K^{\times}(\text{通常拓扑})\to K^{\times}(\text{范数拓扑})
\]
连续, 并在完备化上诱导一个同态 $K^{\times}\to\widehat{K^{\times}}$, 它落在一个交换图表中:
\[
\begin{tikzcd}[column sep=small]
0\ar[r] & U_K\ar[r]\ar[d,"\simeq"] & K^{\times}\ar[r]\ar[d] & \Z\ar[r]\ar[d,hook] & 0\\
0\ar[r] & U_K\ar[r] & \widehat{K^{\times}}\ar[r] & \hat{\Z}\ar[r] & 0
\end{tikzcd}
\]
粗略地说, $\widehat{K^{\times}}$ 是由把 $K^{\times}$ 中的 $\Z$ 换成 $\hat{\Z}$ 而得到的.

\setcounter{theorem}{10} % 1.9, 1.10 为子节标题, 不占定理计数器; 此处将计数器调至 10, 使下一条显示为 1.11
\begin{npara}\label{para:1.11}
选取一个素元 $\pi$ 决定了一个分解
\[
\widehat{K^{\times}}=U_K\cdot\pi^{\hat{\Z}}
\]
为两个闭子群的乘积, 从而 (由无限 Galois 理论) 决定了一个分解
\begin{equation*}
K^{\mathrm{ab}}=K_{\pi}\cdot K^{\mathrm{un}},
\tag{9}
\end{equation*}
其中 $K_{\pi}$ 是 $K^{\mathrm{ab}}$ 中被 $\phi_K(\pi)$ 固定的子域, $K^{\mathrm{un}}$ 是被 $\phi_K(U_K)$ 固定的子域. 显然, $K_{\pi}$ 是全部满足 $\pi\in\Nm(L^{\times})$ 的有限 Abel 扩张 $L/K$ (必全分歧) 的并. 例如, 当在 $K=\Q_p$ 中选取素元 $p$ 时, 分解式成为
\[
\Q_p^{\mathrm{ab}}=\Bigl(\bigcup_n\Q_p[\zeta_{p^n}]\Bigr)\cdot\Bigl(\bigcup_{(m,p)=1}\Q_p[\zeta_m]\Bigr).\erratum{原书右端两处 $\Q$ 均漏了下标 $p$, 依 Milne errata 补正.}
\]
映射 $\phi_K$ 限制为一个同构
\[
U_K\to\Gal(K_{\pi}/K).
\]
这个同构把 \ref{eqn:i.1.9.filt} 式 (本书 p.\pageref{eqn:i.1.9.filt}) 中的滤列映到 $\Gal(K_{\pi}/K)$ 上按分歧群 (上标号; 见 \ref{prop:4.3}) 的滤列.
\end{npara}

\begin{npara}\label{para:1.12}
两个全分歧扩张的复合未必全分歧. 例如考虑 $\Q$ 的二次扩张 $\Q[\sqrt{p}]$ 与 $\Q[\sqrt{pq}]$, 其中 $p$, $q$ 为不同的奇素数. 此时 $p$ 在这两个扩张中都全分歧, 但在它们的复合 $\Q[\sqrt{p},\sqrt{q}]$ 中并非全分歧, 因为它在子域 $\Q[\sqrt{q}]$ 中非分歧. 当 $q$ 取为模 $p$ 的非平方时, 把 $\Q$ 换成 $\Q_p$ 这些陈述仍然成立.

因此, 与 Abel 扩张及非分歧扩张的情形相反, 在 $K^{\mathrm{ab}}$ (或 $K^{\mathrm{al}}$) 中不存在 $K$ 的 “最大” 全分歧扩张: 只有依赖于 $\pi$ 选取的极大全分歧扩张 $K_{\pi}$.
\end{npara}

\subsection*{主定理证明的概述}

我们先说明局部 Artin 映射的唯一性由其余陈述推出.

\begin{theorem}\label{thm:1.13}
假设局部存在定理\ref{thm:1.4}. 至多存在一个同态 $\phi\colon K^{\times}\to\Gal(K^{\mathrm{ab}}/K)$ 满足定理\ref{thm:1.1} 的条件 (a) 与 (b).
\end{theorem}

\begin{proof}
若不存在这样的 $\phi$, 则无可证明; 因此我们假定存在某个 $\phi$, 从而推论\ref{cor:1.2} 以 “有限指数的开子群” 替换 “范数群” 后成立. 设 $\pi$ 是 $\mathcal{O}_K$ 的素元, 令 $K_{\pi,n}$ 是满足 $\Nm(K_{\pi,n})=(1+\mathfrak{m}^n)\cdot\langle\pi\rangle$ 的 $K$ 的扩张, 并令 $K_{\pi}=\bigcup_n K_{\pi,n}$. 由于 $\pi$ 从每个 $K_{\pi,n}$ 都是范数, $\phi(\pi)$ 在 $K_{\pi}$ 上作用为 $1$. 由于它在 $K^{\mathrm{un}}$ 上作用为 Frobenius 自同构, $\phi(\pi)$ 在 $K_{\pi}\cdot K^{\mathrm{un}}$ 上的作用便完全确定. 但 $K_{\pi}\cdot K^{\mathrm{un}}=K^{\mathrm{ab}}$ (参见 \ref{para:1.11}), 因此这表明对任意两个满足条件的同态 $\phi$ 与 $\phi'$ 都有 $\phi(\pi)=\phi'(\pi)$. 由于 $\pi$ 任意, 我们已证明若 $\phi$ 与 $\phi'$ 满足条件, 则对所有 $\pi$ 都有 $\phi(\pi)=\phi'(\pi)$. 群 $K^{\times}$ 由 $\mathcal{O}_K$ 的素元生成, 这就证明了 $\phi=\phi'$.
\end{proof}

剩下要证明存在性, 即局部 Artin 映射与局部存在定理的存在性. 我们将给出三个证明. 前两个的优点在于它们显式地构造了具有范数群 $(1+\mathfrak{m}^n)\cdot\langle\pi\rangle$ 的扩张 $K_{\pi,n}$, 以及局部 Artin 映射 $K^{\times}/\Nm(K_{\pi,n})\to\Gal(K_{\pi,n}/K)$. 后两个的优点在于它们证明了整体类域论中用到的上同调结果. 第三个的优点是简短.

只求最快到达整体类域论主定理的读者, 只需理解第三个证明, 因此可以跳过本章其余部分直接进入第 \ref{ch:II} 章.

\paragraph{存在性的第一个证明 (Lubin-Tate 与 Hasse-Arf; I \textsection\textsection 2--4)}
Lubin 与 Tate 的理论显式地构造了每个素元对应的域 $K_{\pi,n}$, 从而构造了域 $K_{\pi}$, 并且显式地构造了一个同态 $\phi_{\pi}\colon U_K\to\Gal(K_{\pi}/K)$; 此外它还表明 $K_{\pi}\cdot K^{\mathrm{un}}$ 以及 $\phi_{\pi}$ 到同态 $K^{\times}\to\Gal(K_{\pi}\cdot K^{\mathrm{un}}/K)$ (满足 $\phi(\pi)|_{K^{\mathrm{un}}}=\Frob_K$) 的唯一扩张都与 $\pi$ 无关. 由此得到定理\ref{thm:1.1} 与 \ref{cor:1.2}, 但 $K^{\mathrm{ab}}$ 换成了 $K_{\pi}\cdot K^{\mathrm{un}}$.

还需在附注中补入这一点的证明. 令 $K^{\mathrm{LT}}=K_{\pi}\cdot K^{\mathrm{un}}$. 一个要点是, 若 $L$ 是 $K$ 的有限扩张, 则图表
\[
\begin{tikzcd}
L^{\times}\ar[r,"\phi_L"]\ar[d,"\Nm"'] & \Gal(L^{\mathrm{LT}}/L)\ar[d,"\tau\mapsto\tau|_{K^{\mathrm{LT}}}"]\\
K^{\times}\ar[r,"\phi_K"] & \Gal(K^{\mathrm{LT}}/K)
\end{tikzcd}
\]
交换 (\cite{Iwasawa1986}, 定理 6.9, p.\,89). 这表明, 对任何包含在 $K^{\mathrm{LT}}$ 中的 $K$ 的 Abel 扩张 $L$, $\phi_K$ 定义了一个满同态
\[
\phi_{L/K}\colon K^{\times}/\Nm(L^{\times})\to\Gal(L/K).
\]
下一步证明它是单的 (ibid., 推论, p.\,90). 这就证明了局部互反律 (以 $K^{\mathrm{LT}}$ 代 $K^{\mathrm{ab}}$), 而局部存在定理随之易得.

剩下只需证明 $K_{\pi}\cdot K^{\mathrm{un}}=K^{\mathrm{ab}}$. 这一点最直接的证明用到代数数论中的 Hasse-Arf 定理.

\paragraph{主定理的第二个证明 (Lubin-Tate 与上同调; I \textsection\textsection 2--3; III \textsection\textsection 1--3)}
按定理\ref{thm:1.1} 与 \ref{thm:1.4}, 对每个素元 $\pi\in\mathcal{O}_K$ 与整数 $n\ge 1$, 存在 $K$ 的一个 Abel 扩张 $K_{\pi,n}$ 满足
\[
\Nm(K_{\pi,n})=(1+\mathfrak{m}^n)\cdot\pi^{\Z}.
\]
这些扩张具有以下性质:

${}^*\mathrm{a}$ $[K_{\pi,n}:K]=(q-1)q^{n-1}$;

${}^*$b 对所有 $n$, $\pi$ 是 $K_{\pi,n}$ 中的范数.

此外 $K^{\mathrm{ab}}=K_{\pi}\cdot K^{\mathrm{un}}$, 其中 $K_{\pi}=\bigcup_{n\ge 1}K_{\pi,n}$.

在本章 \textsection 3, 我们将对每个素元 $\pi\in\mathcal{O}_K$ 与整数 $n\ge 1$ 显式构造一个满足 ${}^*a$ 与 ${}^*b$ 的全分歧 Abel 扩张 $K_{\pi,n}$. 令 $K_{\pi}=\bigcup_{n\ge 1}K_{\pi,n}$, 并对 $m\ge 1$ 令 $K_m$ 为 $K$ 的 $m$ 次非分歧扩张. 我们将显式构造一个同态
\[
\phi_{\pi}\colon K^{\times}\to\Gal((K_{\pi}\cdot K^{\mathrm{un}})/K)
\]
使得

${}^*\mathrm{c}$ $\phi_{\pi}(\pi)|_{K^{\mathrm{un}}}=\Frob_K$;

${}^*\mathrm{d}$ 对所有 $m$, $n$, 当 $a\in(1+\mathfrak{m}^n)\cdot\langle\pi^m\rangle$ 时 $\phi_{\pi}(a)|_{(K_{\pi,n}\cdot K_m)}=\mathrm{id}$.

此外, 我们将证明 $K_{\pi}\cdot K^{\mathrm{un}}$ 与 $\phi_{\pi}$ 都与 $\pi$ 的选取无关.

在第 \ref{ch:III} 章定理\ref{thm:iii.3.4} 中, 我们将证明存在同态 $\phi\colon K^{\times}\to\Gal(K^{\mathrm{ab}}/K)$ 满足定理\ref{thm:1.1} 的条件 (a) 与 (b).

在本小节的余下部分, 我们说明这两个结果如何推出定理\ref{thm:1.4} (从而也推出定理\ref{thm:1.1}), 并附带更精确的陈述: 对所有 $\pi$ 有 $K^{\mathrm{ab}}=K_{\pi}\cdot K^{\mathrm{un}}$ 且 $\phi=\phi_{\pi}$.

设 $K'=K_{\pi}\cdot K^{\mathrm{un}}$ 是 $K^{\mathrm{ab}}$ 的子域, 并令 $\phi'=\phi_{\pi}$ — 注意 $K'$ 与 $\phi'$ 都与素元 $\pi$ 的选取无关.

\begin{npara}\label{para:1.14}
对所有 $a\in K^{\times}$, $\phi(a)|_{K'}=\phi'(a)$.
\end{npara}

\begin{proof}
对 $K$ 的每个素元 $\pi$, $\phi(\pi)$ 在 $K_{\pi,n}$ 上平凡作用, 因为 $\pi$ 是 $K_{\pi,n}$ 中的范数; 而 $\phi'(\pi)$ 在 $K_{\pi,n}$ 上平凡作用, 是由于条件 ${}^*c$ 中取 $m=1$ (我们可假定定义 $\phi'$ 所用的素元就是 $\pi$). 由于 $\phi(\pi)$ 与 $\phi'(\pi)$ 在 $K^{\mathrm{un}}$ 上都作用为 $\Frob_K$, 它们必须在 $K'=\bigcup_n K_{\pi,n}\cdot K^{\mathrm{un}}$ 上一致. 但 $K$ 的素元作为乘法群生成 $K^{\times}$ (任 $a\in K^{\times}$ 可写成 $a=u\pi^r$, 而 $u=(u\pi)\pi^{-1}$), 故这证明了断言.
\end{proof}

\begin{npara}\label{para:1.15}
\textbf{局部 Kronecker-Weber 定理.} 对任何素元 $\pi$,
\[
K^{\mathrm{ab}}=K_{\pi}\cdot K^{\mathrm{un}}.
\]
\end{npara}

\begin{proof}
令
\[
K_{n,m}=K_{\pi,n}\cdot K_m,\qquad U_{n,m}=(1+\mathfrak{m}^n)\cdot\langle\pi^m\rangle.
\]
已知对所有 $a\in U_{n,m}$ 有 $\phi_{\pi}(a)|_{K_{n,m}}=1$. 因此对所有 $a\in U_{n,m}$ 有 $\phi(a)|_{K_{n,m}}=1$, 从而 $U_{n,m}\subset\Nm(K_{n,m}^{\times})$. 但
\[
\begin{aligned}
(K^{\times}:U_{n,m})&=(U_K:1+\mathfrak{m}^n)(\langle\pi\rangle:\langle\pi^m\rangle)\\
&=(q-1)q^{n-1}\cdot m\\
&=[K_{\pi,n}:K][K_m:K]\\
&=[K_{n,m}:K],
\end{aligned}
\]
且已知 $\phi$ 诱导一个同构
\[
K^{\times}/\Nm(K_{n,m}^{\times})\to\Gal(K_{n,m}/K).
\]
因此
\[
U_{n,m}=\Nm(K_{n,m}^{\times}).
\]
现在设 $L$ 是 $K$ 的有限 Abel 扩张. 已知 $\phi$ 定义同构 $K^{\times}/\Nm(L^{\times})\to\Gal(L/K)$, 故 $\Nm(L^{\times})$ 在 $K^{\times}$ 中指数为有限. 按引理\ref{lem:1.3}, 它也是开的, 从而对某个 $n,m\ge 0$ 包含 $U_{n,m}$. 映射
\[
\phi\colon K^{\times}\to\Gal(L\cdot K_{n,m}/K)
\]
是满的; 且对 $a\in K^{\times}$,
\[
\phi(a)\text{ 固定 }L\text{ 的元素}\iff a\in\Nm(L^{\times}),
\]
\[
\phi(a)\text{ 固定 }K_{n,m}\text{ 的元素}\iff a\in\Nm(K_{n,m}^{\times})=U_{n,m}.
\]
由于 $\Nm(L^{\times})\supset U_{n,m}$, 这推出 $L\subset K_{n,m}$.

由此得到 $K^{\mathrm{ab}}=K_{\pi}\cdot K^{\mathrm{un}}$.

我们现在知道, 对 $K$ 的每个素元 $\pi$, 都有 $K^{\mathrm{ab}}=K_{\pi}\cdot K^{\mathrm{un}}$ 且 $\phi=\phi_{\pi}$. 为完成局部存在定理 (\ref{thm:1.4}) 的证明, 我们必须证明 $K^{\times}$ 的每个有限指数开子群 $H$ 都是范数群; 但如上所观察到的, 每个这样的群都对某个 $n$, $m$ 包含 $U_{n,m}$, 而 $U_{n,m}=\Nm(K_{n,m})$. 于是 (\ref{cor:1.2}e) 表明 $H$ 是范数群.
\end{proof}

\paragraph{主定理的第三个证明 (上同调与 Hilbert 符号; III \textsection\textsection 1--5)}
在第 \ref{ch:III} 章, 我们将用上同调 (具体地, Tate 的一个定理) 来构造局部 Artin 映射, 并用 Hilbert 符号来证明 $K^{\times}$ 的每个有限指数开子群都是范数群.

\begin{exercise}\label{ex:1.16}
只用代数数论的结果 (例如 ANT) 证明: 局部域的有限扩张 $L/K$ 全分歧, 当且仅当 $\Nm(L/K)$ 包含一个素元.
\end{exercise}

\section{Lubin-Tate 形式群律}\label{sec:I.2}

关于 $A$ 上形式群律的一个重要事实是: 它把每个有限扩张 $L/K$ 的整数中的极大理想 $\mathfrak{M}$ 变成一个 Galois 模. 足够小心的话, 你甚至可以让 $L$ 是 $K$ 的一个无限完备扩张. 无论如何, 这个群可能与加法群 $(\mathfrak{M},+)$ 以及乘法群 $(1+\mathfrak{M},\times)$ 都大不相同. 特别地, 它的挠子群可以为 Galois 群提供一个非常有用的表示模. (Lubin, \sx{2061570})

\subsection*{幂级数}

设 $A$ 是一个环 (总假定交换且带 $1$). 一个系数在 $A$ 中的幂级数是一个无穷序列
\[
f=(a_0,a_1,a_2,\ldots),\quad a_i\in A,\ i\in\N.
\]
加法与乘法定义为
\[
(a_0,a_1,\ldots)+(b_0,b_1,\ldots)=(a_0+b_0,a_1+b_1,\ldots),
\]
\[
(a_0,a_1,\ldots)(b_0,b_1,\ldots)=(a_0b_0,\ldots,\textstyle\sum_{i+j=k}a_i b_j,\ldots).
\]
若写成
\[
f=\sum_{i\ge 0}a_i T^i,
\]
这些公式便易于记忆.

系数在 $A$ 中的幂级数构成一个交换环, 记为 $A[[T]]$. 幂级数可以像多项式一样处理, 但有几点需要注意. 例如, 一般我们不能把元素 $c\in A$ 代入幂级数 $f(T)\in A[[T]]$, 因为计算 $f(c)=\sum_{i\ge 0}a_i c^i$ 要求我们求和 $A$ 中无穷多个元素, 而我们不是分析学家, 做不到这一点. 同样的理由, 我们只能在一个幂级数 $g(T)$ 的常数项为零时把它代入另一个幂级数 $f(T)$; 此时 $f(g(T))$ 有定义, 记为 $f\circ g$.

\begin{lemma}\label{lem:2.1}
(a) 对所有幂级数 $f\in A[[T]]$ 与 $g,h\in TA[[T]]$,
\[
f\circ(g\circ h)=(f\circ g)\circ h.
\]
(b) 设 $f=\sum_{i\ge 1}^{\infty}a_i T^i\in TA[[T]]$. 存在 $g\in TA[[T]]$ 使得 $f\circ g=T$, 当且仅当 $a_1$ 是 $A$ 中的单位; 此时 $g$ 唯一, 且满足 $g\circ f=T$.
\end{lemma}

\begin{proof}
(a) 一般有 $(f_1f_2)\circ g=(f_1\circ g)(f_2\circ g)$, 从而 $f^n\circ g=(f\circ g)^n$. 因此当 $f=T^n$ 时, $f\circ(g\circ h)$ 与 $(f\circ g)\circ h$ 都等于 $(g\circ h)^n$; 当 $f=\sum a_i T^i$ 时, 两者都等于 $\sum a_i(g\circ h)^i$.

(b) 我们寻找 $g=\sum_{i\ge 1}b_i T^i$ 使得
\[
\sum_{i\ge 1}a_i g^i=T,
\]
即
\[
a_1b_1=1,
\]
\[
a_1b_2+a_2b_1^2=0,
\]
\[
a_1b_n+\bigl(\text{关于 }a_2,\ldots,a_n,b_1,\ldots,b_{n-1}\text{ 的多项式}\bigr)=0.
\]
第一个方程表明, 为使 $g$ 存在, $a_1$ 必须可逆. 当 $a_1$ 可逆时, 这些方程唯一地确定了 $b_i$. 现在, 由于 $b_1$ 可逆, 同样的论证表明存在 $h\in TA[[T]]$ 使得 $g\circ h=T$. 但
\[
f=f\circ T=f\circ g\circ h=T\circ h=h,
\]
故 $g\circ f=T$.
\end{proof}

注意: 一般有 $f\circ(g+h)\neq f\circ g+f\circ h$.

多个变元的幂级数可类似定义. 此外, 若 $f(X_1,\ldots,X_n)\in A[[X_1,\ldots,X_n]]$ 且 $g_1,g_2,\ldots,g_n\in A[[Y_1,\ldots,Y_m]]$, 则只要所有 $g_i$ 的常数项都为零, $f(g_1,\ldots,g_n)$ 就是 $A[[Y_1,\ldots,Y_m]]$ 中一个良定义的元素.

\begin{remark}\label{rem:2.2}
设 $A$ 是一个完备离散赋值环, $\mathfrak{m}$ 是 $A$ 中的极大理想. 对每个 $f=\sum_{i\ge 0}a_i T^i\in A[[T]]$ 与每个 $c\in\mathfrak{m}$, 当 $i\to\infty$ 时 $a_i c^i\to 0$. 因此级数 $\sum_{i\ge 0}a_i c^i$ 收敛到 $A$ 中的一个元素 $f(c)$, 且当 $a_0\in\mathfrak{m}$ 时它落在 $\mathfrak{m}$ 中.
\end{remark}

\subsection*{形式群律}

一个群是一个带有满足群公理的合成法则的非空集合. 一个\emph{形式群律}\index{xingshiqunlü@形式群律(formal group law)}是一个满足群公理的合成法则 (但没有集合). 更精确地:

\begin{definition}\label{def:2.3}
设 $A$ 是一个交换环. $A$ 上的一个\emph{一元交换形式群律}是一个幂级数 $F\in A[[X,Y]]$ 满足

(a) $F(X,Y)=X+Y+\text{次数}\ge 2\text{ 的项}$;

(b) $F(X,F(Y,Z))=F(F(X,Y),Z)$;

(c) 存在唯一的 $i_F(X)\in XA[[X]]$ 使得 $F(X,i_F(X))=0$;

(d) $F(X,Y)=F(Y,X)$.
\end{definition}

\begin{remark}\label{rem:2.4}
(a) 条件 (a) 保证 $F(X,Y)$ 无常数项, 因此公理 (b) 有意义: 我们在每个次数上比较有限和.

(b) 在公理 (a) 与 (b) 中取 $Y=Z=0$, 我们发现
\[
F(X,0)=X+\text{次数}\ge 2\text{ 的项},\qquad F(F(X,0),0)=F(X,0).
\]
把幂级数 $F(X,0)$ 记为 $f(X)$. 第一个等式表明存在 $g$ 使得 $f\circ g=X$, 第二个等式说 $f\circ f=f$. 在第二个等式上复合 $g$ 得 $f=X$. 于是 $F(X,0)=X$, 类似地 $F(0,Y)=Y$. 因此
\[
F(X,Y)=X+Y+\sum_{\substack{1\le i<\infty\\ 1\le j<\infty}}a_{i,j}X^iY^j.
\]
要得到 $n$ 元群律, 把 $X$ 与 $Y$ 各自换成 $n$ 个变元的序列即可. 公理 (d) 是交换性条件. 由于我们不再考虑其他情形, 以后简称一元交换形式群律为形式群律.
\end{remark}

设 $A=\mathcal{O}_K$ 是非阿基米德局部域 $K$ 中的整数环, 并设 $F=\sum_{i,j}a_{ij}X^iY^j$ 是 $\mathcal{O}_K$ 上的形式群律. 对每个 $x,y\in\mathfrak{m}_K$, 当 $(i,j)\to\infty$ 时 $a_{ij}x^i y^j\to 0$, 因此级数
\[
F(x,y)=\sum a_{ij}x^i y^j
\]
收敛到 $\mathfrak{m}_K$ 中的一个元素 $x+_F y$. 这样 $\mathfrak{m}_K$ 成为一个交换群 $(\mathfrak{m}_K,+_F)$. 类似地, 对每个有限扩张 $L\supset K$, $\mathfrak{m}_L$ 获得一个群结构, 且包含映射 $(\mathfrak{m}_K,+_F)\hookrightarrow(\mathfrak{m}_L,+_F)$ 是同态.

\begin{example}\label{ex:2.5}
(a) 设 $F(X,Y)=X+Y$. 则 $+_F$ 是 $\mathfrak{m}_K$ 上的通常加法.

(b) 设 $F(X,Y)=X+Y+XY$. 映射
\[
a\mapsto 1+a\colon\mathfrak{m}\to 1+\mathfrak{m}
\]
是一个同构 $(\mathfrak{m},+_F)\to(1+\mathfrak{m},\times)$. 验证:
\[
\begin{tikzcd}
(a,b)\ar[r,mapsto]\ar[d,"+_F"'] & (1+a,1+b)\ar[d,"\times"]\\
a+b+ab\ar[r,mapsto] & 1+a+b+ab.
\end{tikzcd}
\]

(c) 设 $E$ 是非阿基米德局部域 $K$ 上的一条椭圆曲线, 并设
\[
Y^2Z+a_1XYZ+a_3YZ^2=X^3+a_2X^2Z+a_4XZ^2+a_6Z^3
\]
是 $E$ 的一个极小 Weierstrass 模型. 令 $T=-X/Y$. 把 $E$ 上的群律展开为 $T_1$, $T_2$ 的幂级数, 我们得到 $\mathcal{O}_K$ 上的一个形式群律 $F_E(T_1,T_2)$. 见 \cite{Milne2020}, 2.6, 2.7.
\end{example}

\begin{definition}\label{def:2.6}
设 $F(X,Y)$ 与 $G(X,Y)$ 是形式群律. 一个\emph{同态}\index{xingshiqunlü@形式群律(formal group law)!tonghui@同态(homomorphism)} $F\to G$ 是一个幂级数 $h\in TA[[T]]$ 使得
\[
h(F(X,Y))=G(h(X),h(Y)).
\]
当 $h$ 有逆, 即存在同态 $h'\colon G\to F$ 使得
\[
h\circ h'=T=h'\circ h
\]
时, $h$ 称为\emph{同构}\index{xingshiqunlü@形式群律(formal group law)!tonggou@同构(isomorphism)}. 一个同态 $h\colon F\to F$ 称为 $F$ 的\emph{自同态}\index{xingshiqunlü@形式群律(formal group law)!zitongtai@自同态(endomorphism)}.
\end{definition}

在 $A=\mathcal{O}_K$ 的情形, 一个同态 $f\colon F\to G$ 对每个 $L\supset K$ 定义一个同态
\[
a\mapsto f(a)\colon(\mathfrak{m}_L,+_F)\to(\mathfrak{m}_L,+_G).
\]

\begin{example}\label{ex:2.7}
设 $F=X+Y+XY=(1+X)(1+Y)-1$. 则 $f(T)=(1+T)^p-1$ 是 $F$ 的一个自同态, 因为
\[
F(f(X),f(Y))=(1+X)^p(1+Y)^p-1=f(F(X,Y)).
\]
注意下列图表交换:
\[
\begin{tikzcd}
\mathfrak{m}\ar[r,"f"]\ar[d,"a\mapsto 1+a"'] & \mathfrak{m}\ar[d,"a\mapsto 1+a"]\\
1+\mathfrak{m}\ar[r,"a\mapsto a^p"] & 1+\mathfrak{m}
\end{tikzcd}
\]
即当把 $(\mathfrak{m},+_F)$ 与 $(1+\mathfrak{m},\times)$ 等同起来时, $f$ 就变成了 $a\mapsto a^p$.
\end{example}

设 $G$ 是一个形式群律. 对任意 $f,g\in TA[[T]]$, 我们定义
\[
f+_G g=G(f(T),g(T)).
\]
由公理 \ref{def:2.3}a,b,c,d, 这个合成法则使 $TA[[T]]$ 成为一个交换群. 特别地,
\[
f+_G(i_G\circ f)=0.
\]

\begin{lemma}\label{lem:2.8}
(a) 对任何形式群律 $F$ 与 $G$, 从 $F$ 到 $G$ 的同态的集合 $\Hom(F,G)$ 关于加法 $f+_G g$ 成为一个 Abel 群.

(b) 对任何形式群律 $F$, $F$ 的自同态的 Abel 群 $\End(F)$ 关于乘法 $f\circ g$ 成为一个环 (未必交换).
\end{lemma}

\begin{proof}
设 $f$ 与 $g$ 是 $F\to G$ 的同态, 并令 $h=f+_G g$. 则
\[
\begin{aligned}
h(F(X,Y))&\stackrel{\mathrm{def}}{=}G(f(F(X,Y)),g(F(X,Y)))\\
&=G(G(f(X),f(Y)),G(g(X),g(Y))).
\end{aligned}
\]
(至少在符号上) 可把这个最后的幂级数写成
\begin{equation*}
(f(X)+_G f(Y))+_G(g(X)+_G g(Y)),\tag{$*$}
\end{equation*}
而由结合律与交换律我们可把它改写为
\begin{equation*}
(f(X)+_G g(X))+_G(f(Y)+_G g(Y)),\tag{$**$}
\end{equation*}
即为 $G(h(X),h(Y))$. 更形式地说, 把 (*) 变成 (**) 的那些运算, 也把 $G(G(f(X),f(Y)),G(g(X),g(Y)))$ 变成 $G(h(X),h(Y))$. 这就证明了 $h\in\Hom(F,G)$. 类似地可证 $i_G\circ f\in\Hom(F,G)$. 由于 $0\in\Hom(F,G)$, 这完成了 $\Hom(F,G)$ 是 $(TA[[T]],+_G)$ 的子群的证明.

我们在引理\ref{lem:2.1} 中已证 $f,g\mapsto f\circ g$ 是结合的. 为证 $\End(F)$ 是环, 还需注意对 $f,g,h\in\End(F)$,
\[
\begin{aligned}
f\circ(g+_F h)&\stackrel{\mathrm{def}}{=}f(F(g(T),h(T)))=F((f\circ g)(T),(f\circ h)(T))=f\circ g+_F f\circ h,\\
(f+_F g)\circ h&=f\circ h+_F g\circ h\quad(\text{同理}),
\end{aligned}
\]
且 $\End(F)$ 有单位元, 即 $T$.
\end{proof}

形式群律类似于代数群. 主要差别在于: 由于它们由幂级数而非多项式定义, 其点的坐标必须 “接近 $1$” 才能使乘积有定义. 关于形式群律有一个非常广的理论 — 例如见 M.\,Hazewinkel, \emph{Formal Groups and Applications}, Academic Press, 1978.

\subsection*{Lubin-Tate 群律}

现在令 $A=\mathcal{O}_K$ 为非阿基米德局部域 $K$ 的整数环, 并选取 $A$ 的一个素元 $\pi$.

\begin{definition}\label{def:2.9}
设 $\mathcal{F}_{\pi}$ 是满足下列条件的 $f(X)\in A[[X]]$ 的集合:

(a) $f(X)=\pi X+\text{次数}\ge 2\text{ 的项}$;

(b) $f(X)\equiv X^q\bmod\pi$.
\end{definition}

\begin{example}\label{ex:2.10}
(a) 多项式 $f(X)=\pi X+X^q$ 属于 $\mathcal{F}_{\pi}$.

(b) 设 $K=\Q_p$, $\pi=p$; 则多项式
\[
f(X)=(1+X)^p-1=pX+\binom{p}{2}X^2+\cdots+pX^{p-1}+X^p
\]
属于 $\mathcal{F}_p$.
\end{example}

\begin{lemma}\label{lem:2.11}
设 $f,g\in\mathcal{F}_{\pi}$, 并设 $\phi_1(X_1,\ldots,X_n)$ 是一个系数在 $A$ 中的线性型. 存在唯一的 $\phi\in A[[X_1,\ldots,X_n]]$ 使得
\[
\begin{cases}
\phi(X_1,\ldots,X_n)=\phi_1+\text{次数}\ge 2\text{ 的项},\\
f(\phi(X_1,\ldots,X_n))=\phi(g(X_1),\ldots,g(X_n)).
\end{cases}
\]
\end{lemma}

\begin{proof}
我们对 $r$ 用归纳法证明: 存在唯一的 $r$ 次多项式 $\phi_r(X_1,\ldots,X_n)$ 使得
\[
\begin{cases}
\phi_r(X_1,\ldots,X_n)=\phi_1+\text{次数}\ge 2\text{ 的项},\\
f(\phi_r(X_1,\ldots,X_n))=\phi_r(g(X_1),\ldots,g(X_n))+\text{次数}\ge r+1\text{ 的项}.
\end{cases}
\]
第一个多项式的唯一候选就是 $\phi_1$ 本身. 它显然满足第一个条件; 若写 $\phi_1=\sum a_i X_i$, 则第二个条件说
\[
\pi\Bigl(\sum a_i X_i\Bigr)=\sum a_i(\pi X_i)+\text{次数}\ge 2\text{ 的项},
\]
这也成立.

设 $r\ge 1$ 且我们已定义 $\phi_r$. 由于 $\phi_r$ 唯一, $\phi_{r+1}$ 必等于 $\phi_r+Q$, 其中 $Q$ 是 $A[X_1,\ldots,X_n]$ 中一个 $r+1$ 次齐次多项式. 我们需要
\[
f(\phi_{r+1}(X_1,\ldots,X_n))\stackrel{?}{=}\phi_{r+1}(g(X_1),\ldots,g(X_n))+\text{次数}\ge r+2\text{ 的项}.
\]
左边为
\[
f(\phi_r(X_1,\ldots,X_n))+\pi Q(X_1,\ldots,X_n)+\text{次数}\ge r+2\text{ 的项},
\]
而右边为
\[
\phi_r(g(X_1),\ldots,g(X_n))+Q(\pi X_1,\ldots,\pi X_n)+\text{次数}\ge r+2\text{ 的项}.
\]
由于 $Q$ 是 $r+1$ 次齐次的, $Q(\pi X_1,\ldots)=\pi^{r+1}Q(X_1,\ldots)$,\erratum{原书右端误作 $\pi^{r+1}Q(\pi X_1,\ldots)$, 多出一个 $\pi$, 依 Milne errata (D. Craven) 删去.} 故我们需要
\[
(\pi^{r+1}-\pi)Q(X_1,\ldots,X_n)\stackrel{?}{=}f(\phi_r(X_1,\ldots,X_n))-\phi_r(g(X_1),\ldots,g(X_n))+\text{次数}\ge r+2\text{ 的项}.
\]
因此 $Q$ 必须是满足下式的唯一多项式:
\[
\frac{f(\phi_r(X_1,\ldots,X_n))-\phi_r(g(X_1),\ldots,g(X_n))}{(\pi^r-1)\pi}=Q+\text{次数}\ge r+2\text{ 的项}.
\]
注意, 由于二项式定理在特征 $p$ 下的简单形式,
\[
f\circ\phi_r-\phi_r\circ g\equiv\phi_r(X_1,\ldots,X_n)^q-\phi_r(X_1^q,\ldots,X_n^q)\equiv 0\bmod\pi.
\]
由于 $\pi$ 整除 $f\circ\phi_r-\phi_r\circ g$, 而 $\pi^r-1$ 在 $A$ 中可逆, $Q$ 的系数确实在 $A$ 中; 又因 $\phi_r$ 满足归纳假设, 它的次数确为 $r+1$.

定义了 $\phi_r$ ($r=1,2,\ldots$) 并注意
\[
\phi_{r+1}=\phi_r+\text{次数}\ge r+1\text{ 的项}
\]
之后, 我们可以定义 $\phi$ 为满足
\[
\phi=\phi_r+\text{次数}\ge r+1\text{ 的项}\quad(\text{对所有 }r)
\]
的唯一幂级数. 显然它具有所需的第一个性质; 且对任意 $r$,
\[
\begin{aligned}
f(\phi(X_1,\ldots,X_n))&=f(\phi_r(X_1,\ldots,X_n))+\text{次数}\ge r+1\text{ 的项}\\
&=\phi_r(g(X_1,\ldots,X_n))+\text{次数}\ge r+1\text{ 的项}\\
&=\phi(g(X_1,\ldots,X_n))+\text{次数}\ge r+1\text{ 的项}.
\end{aligned}
\]
由于这对所有 $r$ 成立, $\phi$ 也具有所需的第二个性质.
\end{proof}

\begin{proposition}\label{prop:2.12}
对每个 $f\in\mathcal{F}_{\pi}$, 存在唯一的系数在 $A$ 中的形式群律 $F_f$ 以 $f$ 为自同态.
\end{proposition}

\begin{proof}
按引理\ref{lem:2.11}, 存在唯一的幂级数 $F_f(X,Y)$ 使得
\[
\begin{cases}
F_f(X,Y)=X+Y+\text{次数}\ge 2\text{ 的项},\\
f(F_f(X,Y))=F_f(f(X),f(Y)).
\end{cases}
\]
剩下只需验证这是一个形式群律.

交换性: 令 $G=F_f(Y,X)$. 则
\[
\begin{cases}
G(X,Y)=X+Y+\text{次数}\ge 2\text{ 的项},\\
f(G(X,Y))=f(F_f(Y,X))=F_f(f(Y),f(X))=G(f(X),f(Y)).
\end{cases}
\]
由于 $F_f(X,Y)$ 是具有这些性质的唯一幂级数, 推出 $G(X,Y)=F_f(X,Y)$.

结合性: 令 $G_1(X,Y,Z)=F_f(X,F_f(Y,Z))$, $G_2(X,Y,Z)=F_f(F_f(X,Y),Z)$. 则对 $i=1,2$,
\[
\begin{cases}
G_i(X,Y,Z)=X+Y+Z+\text{次数}\ge 2\text{ 的项},\\
G_i(f(X),f(Y),f(Z))=f(G_i(X,Y,Z)),
\end{cases}
\]
而引理\ref{lem:2.11} 再次表明满足这些条件的幂级数只有一个.
\end{proof}

\begin{example}\label{ex:2.13}
设 $K=\Q_p$, $\pi=p$. 则 $f(T)=(1+T)^p-1$ 属于 $\mathcal{F}_p$, 且 $F(X,Y)=X+Y+XY$ 以 $f$ 为自同态 (见 \ref{ex:2.7}). 因此 $F=F_f$.
\end{example}

命题所定义的那些形式群律 $F_f$ 就是 \emph{Lubin-Tate 形式群律}\index{lubintatexingshiqunlü@Lubin-Tate 形式群律(Lubin-Tate formal group)}. 它们恰是以如下方式自同态的形式群律:

\begin{itemize}
\item 其在原点的导数等于 $K$ 的一个素元, 且
\item 其模 $\mathfrak{m}$ 约化到 Frobenius 映射 $T\mapsto T^q$.
\end{itemize}

\begin{proposition}\label{prop:2.14}
对 $f,g\in\mathcal{F}_{\pi}$ 与 $a\in A$, 令 $[a]_{g,f}$ 为 $A[[T]]$ 中满足
\[
\begin{cases}
[a]_{g,f}(T)=aT+\text{次数}\ge 2\text{ 的项},\\
g\circ[a]_{g,f}=[a]_{g,f}\circ f
\end{cases}
\]
的唯一元素. 则 $[a]_{g,f}$ 是一个同态 $F_f\to F_g$.
\end{proposition}

\begin{proof}
令 $h=[a]_{g,f}$ — 它的存在性由引理\ref{lem:2.11} 保证. 我们要证明
\[
h(F_f(X,Y))=F_g(h(X),h(Y)).
\]
显然两边都等于 $aX+aY+\text{次数}\ge 2\text{ 的项}$. 此外,
\[
h(F_f(f(X),f(Y)))=(h\circ f)(F_f(X,Y))=g(h(F_f(X,Y))),
\]
\[
F_g(h(f(X)),h(f(Y)))=F_g(g(h(X)),g(h(Y)))=g(F_g(h(X),h(Y))),
\]
我们可以再次应用引理\ref{lem:2.11} 的唯一性.
\end{proof}

\begin{proposition}\label{prop:2.15}
对任意 $a,b\in A$,
\[
[a+b]_{g,f}=[a]_{g,f}+_{F_g}[b]_{g,f}
\]
且
\[
[ab]_{h,f}=[a]_{h,g}\circ[b]_{g,f}.
\]
\end{proposition}

\begin{proof}
在每种情形, 右边的幂级数都满足刻画左边幂级数的条件.
\end{proof}

\begin{corollary}\label{cor:2.16}
对 $f,g\in\mathcal{F}_{\pi}$, $F_f\simeq F_g$.
\end{corollary}

\begin{proof}
对每个 $u\in A^{\times}$, $[u]_{f,g}$ 与 $[u^{-1}]_{g,f}$ 是互逆的同构.

事实上, 存在唯一的同构 $h\colon F_f\to F_g$ 使得 $h(T)=T+\cdots$ 且 $g\circ h=h\circ f$, 即 $[1]_{g,f}$. 因此 $F_f\simeq F_g$.
\end{proof}

\begin{corollary}\label{cor:2.17}
对每个 $a\in A$, 存在唯一的自同态 $[a]_f\colon F_f\to F_f$ 使得 $[a]_f=aT+\text{次数}\ge 2\text{ 的项}$ 且 $[a]_f$ 与 $f$ 交换. 映射
\[
a\mapsto[a]_f\colon A\hookrightarrow\End(F_f)
\]
是一个环同态.
\end{corollary}

\begin{proof}
取 $[a]_f=[a]_{f,f}$ — 它是满足 $aT+\cdots$ 且与 $f$ 交换的唯一幂级数, 并且是 $F_f$ 的自同态. $a\mapsto[a]_f$ 是环同态这一事实由引理\ref{prop:2.15} 以及显然的事实 $[1]_f=T$ 推出.
\end{proof}

由此可知, 对 $K$ 的有限扩张 $L$, Abel 群 $(\mathfrak{m}_L,+_{F_f})$ 有一个自然的 $A$-模结构.

\begin{example}\label{ex:2.18}
设 $K=\Q_p$, $f=(1+T)^p-1\in\mathcal{F}_p$, 于是 $F_f=X+Y+XY$. 对每个 $a\in\Z_p$, 定义
\[
(1+T)^a=\sum_{m\ge 0}\binom{a}{m}T^m,\qquad \binom{a}{m}\stackrel{\mathrm{def}}{=}\frac{a(a-1)\cdots(a-m+1)}{m(m-1)\cdots 1}.
\]
当 $a\in\Z$ 时, 这些定义与通常的相符; 若 $(a_i)_{i\ge 1}$ 是收敛到 $a\in\Z_p$ 的整数序列, 则当 $i\to\infty$ 时 $\binom{a_i}{m}\to\binom{a}{m}$. 因此 $\binom{a}{m}\in\Z_p$. 我断言
\[
[a]_f=(1+T)^a-1.
\]
显然 $(1+T)^a-1=aT+\cdots$, 且
\[
f\circ((1+T)^a-1)=(1+T)^{ap}-1=((1+T)^a-1)\circ f
\]
当 $a$ 为整数时成立, 由连续性它对所有 $a\in\Z_p$ 成立.

在同构 $(\mathfrak{m},+_{F_f})\xrightarrow{t\mapsto 1+t}(1+\mathfrak{m},\times)$ 下, $[a]_f$ 的作用对应于把 $1+\mathfrak{m}$ 中元素映到其 $a$ 次幂的映射.
\end{example}

\begin{remark}\label{rem:2.19}
(a) 注意 $[\pi]_f=f$, 因为 $f$ 满足那两个定义条件.

(b) 同态 $a\mapsto[a]_f\colon A\to\End(F_f)$ 是单的, 因为 $a$ 可作为 $[a]_f$ 的首项系数恢复出来.

(c) 标准同构 $[1]_{g,f}\colon F_f\to F_g$ 与 $A$ 在 $F_f$ 和 $F_g$ 上的作用交换, 因为
\[
[a]_g\circ[1]_{g,f}=[a]_{g,f}=[1]_{g,f}\circ[a]_f.
\]
\end{remark}

\begin{summary}\label{note:2.20}
对每个 $f\in\mathcal{F}_{\pi}$, 存在唯一一个以 $f$ 为自同态的形式群律 $F_f$. 此外, 在 $F_f$ 上存在一个唯一的 $A$-模结构 $a\mapsto[a]_f\colon A\to\End(F_f)$ 使得

(a) 对所有 $a\in A$, $[a]_f=aT+\text{次数}\ge 2\text{ 的项}$;

(b) $[a]_f$ 与 $f$ 交换.

我们有 $[\pi]_f=f$. 若 $g\in\mathcal{F}_{\pi}$, 则 $F_f\simeq F_g$ (通过一个典范的 $A$-同构).
\end{summary}

\begin{exercise}\label{ex:2.21}
设 $F(X,Y)$ 是一个满足 $F(X,0)=X$ 且 $F(0,Y)=Y$ 的幂级数. 证明存在唯一的幂级数 $G(X)=-X+\sum_{i=2}^{\infty}a_i X^i$ 使得 $F(X,G(X))=0$. 由此推出定义\ref{def:2.3} 中的公理 (c) 是多余的.
\end{exercise}

\begin{NOTES}\label{note:2.22}
设 $f\in\mathcal{F}_{\pi}$. 存在唯一的幂级数 $\log(T)=T+\cdots\in K[[T]]$ 使得 $\log(F(x,y))=\log(x)+\log(y)$. 它具有性质
\[
\log([a]_f(T))=a\log(T)
\]
对 $a\in A$ 成立. 见 \sx{3644403}.
\end{NOTES}

\section{扩张 \texorpdfstring{$K_{\pi}$}{K pi} 的构造}\label{sec:I.3}

在本节, $A=\mathcal{O}_K$, 其中 $K$ 是非阿基米德局部域, 其剩余类域 $A/\mathfrak{m}=k$ 有 $q$ (一个 $p$ 的幂) 个元素. 我们固定 $A$ 的一个素元 $\pi$.

按 \textsection 1 中的讨论, 在 $K^{\mathrm{al}}$ 中应存在唯一的 $K$ 的扩张 $K_{\pi}$ 使得 $K^{\mathrm{ab}}=K_{\pi}\cdot K^{\mathrm{un}}$, 且 $\pi$ 是 $K_{\pi}$ 的每个有限子扩张中的范数, 即 $K^{\mathrm{ab}}$ 中被 $\phi(\pi)$ 固定的子域 $K_{\pi}$.

构造 $K^{\mathrm{un}}$ 很容易. 令 $\mu_m$ 为 $K^{\mathrm{al}}$ 中 $m$ 次单位根的集合, 即 $X^m-1$ 的根的集合. 当 $m$ 不被 $p$ 整除时, $X^m-1$ 的判别式是 $\mathcal{O}_K$ 中的单位, 因此由 $\mu_m$ 中元素生成的域 $K[\mu_m]$ 在 $K$ 上非分歧; 此外 $K[\mu_m]$ 的剩余类域是 $X^m-1$ 在 $k$ 上的分裂域, 故有 $q^f$ 个元素, 其中 $f$ 是使得 $m\mid p^f-1$ 的最小正整数. 因此 $\bigcup_{p\nmid m}K[\mu_m]$ 是 $K$ 的一个非分歧扩张, 其剩余类域为 $k$ 的代数闭包, 从而等于 $K^{\mathrm{un}}$. Galois 群 $\Gal(K^{\mathrm{un}}/K)\simeq\hat{\Z}$, 且 $a\in\hat{\Z}$ 按如下方式作用 $K^{\mathrm{un}}$: 对每个 $\zeta\in\mu_m$ 以及每个充分接近 $a$ 的整数 $a_0$ (依赖于 $m$), $a*\zeta=\zeta^{q^{a_0}}$ ($=\Frob_K^{a_0}(\zeta)$).\erratum{原书误作 $\zeta^{a_0}$ ($=\Frob_K^{a}(\zeta)$), 依 Milne errata (R. Sundararajan) 改正.}

在 $K=\Q_p$, $\pi=p$ 的情形, 对 $K_{\pi}$ 有一个类似的构造, 即 $(\Q_p)_p=\bigcup\Q_p[\mu_{p^n}]$ — 我们以后将证明它具有所指出的性质. 作用
\[
([m],\zeta)\mapsto\zeta^m\colon\Z/p^n\Z\times\mu_{p^n}\to\mu_{p^n}
\]
使 $\mu_{p^n}$ 成为一个秩 $1$ 的自由 $\Z/p^n\Z$-模. 由于 $\Z/p^n\Z\simeq\Z_p/p^n\Z_p$, 我们可以把 $\mu_{p^n}$ 看成一个 $\Z_p$-模, 同构于 $\Z_p/(p^n)$. $\Z_p$ 在 $\mu_{p^n}$ 上的作用诱导一个同构 $(\Z_p/p^n\Z_p)^{\times}\to\Gal(\Q_p[\mu_{p^n}]/\Q_p)$, 取所有 $n$ 的极限得到同构
\[
\Z_p^{\times}\to\Gal((\Q_p)_p/\Q_p).
\]
因此, 对扩张 $K^{\mathrm{un}}/K$ 与 $(\Q_p)_p/\Q_p$ 两者, 我们都有扩张的一组显式生成元, Galois 群的一个显式描述, 以及 Galois 群在生成元集合上的一个显式描述. 引人注目的是, Lubin-Tate 群对每个 $K$ 与 $\pi$ 的 $K_{\pi}/K$ 给出了类似的结果.

$K$ 上的绝对值 $|\cdot|$ 唯一地扩张到每个有限次子域 $L\subset K^{\mathrm{al}}$, 从而扩张到 $K^{\mathrm{al}}$; 将 $K^{\mathrm{al}}$ 对 $|\cdot|$ 完备化, 记所得的完备域为 $\widehat{K^{\mathrm{al}}}$. 设 $f\in\mathcal{F}_{\pi}$. 对任何 $\alpha,\beta\in\widehat{K^{\mathrm{al}}}$ 满足 $|\alpha|,|\beta|<1$ 以及 $a\in A$, 级数 $F_f(\alpha,\beta)$ 与 $[a]_f(\alpha)$ 收敛.\erratum{原书直接在 $K^{\mathrm{al}}$ 中断言这些级数收敛; 但 $K^{\mathrm{al}}$ 不完备, 收敛性并无保证. 依 Milne errata (R. Sundararajan), 改在其完备化 $\widehat{K^{\mathrm{al}}}$ 中陈述; 各 $\pi^n$-挠子模 $\Lambda_n$ 作为 $f^{(n)}$ 的根之集仍含于 $K^{\mathrm{al}}$ (见 \ref{rem:3.1}).} 因此我们可以定义 $A$-模 $\Lambda_f$ 为:
\[
\begin{aligned}
\Lambda_f&=\{\alpha\in\widehat{K^{\mathrm{al}}}\mid |\alpha|<1\}\quad(\text{作为集合}),\\
\alpha+_{\Lambda_f}\beta&=\alpha+_F\beta=F_f(\alpha,\beta),\\
a*\alpha&=[a]_f(\alpha).
\end{aligned}
\]
我们定义 $\Lambda_n$ 为 $\Lambda_f$ 中被 $[\pi]_f^n$ 零化的元素的子模. 若 $g$ 是 $\mathcal{F}_{\pi}$ 中另一个元素, 则典范 $A$-同构 $F_f\to F_g$ 诱导一个 $A$-同构 $\Lambda_f\to\Lambda_g$.

\begin{remark}\label{rem:3.1}
回忆 $[\pi]_f(T)=f(T)$, 因此 $\Lambda_n$ 是
\[
f^{(n)}\stackrel{\mathrm{def}}{=}f\circ f\circ\cdots\circ f\qquad(n\text{ 个因子})
\]
在 $K^{\mathrm{al}}$ 中绝对值 $<1$ 的根的集合. 为简单起见, 我们取 $f$ 为多项式 $\pi T+a_2T^2+\cdots+T^q$ 而非更复杂的幂级数 — 按 \ref{cor:2.16} 甚至取 $f=\pi T+T^q$ 就足够.\footnote{或者, $p$-进 Weierstrass 预备定理 (\cite{Washington1997}, 7.3) 表明每个 $f\in F_{\pi}$ 都可分解为 $f_{1}(T)u(T)$, 其中 $f_{1}(T)$ 是多项式 $\pi T+\cdots+aT^{q}$ ($a\equiv1\mod\mathfrak{m}$), 而 $u(T)$ 是 $A[[T]]$ 中的单位. 显然 $f^{(n)}$ 的根就是 $f_{1}^{(n)}$ 的根。} 于是
\[
\begin{aligned}
(f\circ f)(T)&\stackrel{\mathrm{def}}{=}f(f(T))=\pi(\pi T+\cdots+T^q)+\cdots+(\pi T+\cdots+T^q)^q\\
&=\pi^2 T+\cdots+T^{q^2},
\end{aligned}
\]
且
\[
f^{(n)}(T)=\pi^n T+\cdots+T^{q^n}.
\]
由 $f^{(n)}$ 的 Newton 多边形 (参见 ANT, 7.44) 我们看到 $f^{(n)}$ 的根都有正的 $\mathrm{ord}_K$, 从而绝对值 $<1$, 因此 $\Lambda_n$ 是 $f^{(n)}$ 在 $K^{\mathrm{al}}$ 中的全部根, 赋予交换群结构
\[
\alpha+_{F_f}\beta=F_f(\alpha,\beta)=\alpha+\beta+\cdots
\]
以及 $A$-模结构
\[
[a]_f\alpha=a\alpha+\cdots.
\]
\end{remark}

\begin{example}\label{ex:3.2}
取 $K=\Q_p$, $f=(T+1)^p-1\in\mathcal{F}_p$; 则
\[
\Lambda_n=\{\alpha\in\Q_p^{\mathrm{al}}\mid(\alpha+1)^{p^n}=1\}\simeq\{\zeta\mid\zeta^{p^n}=1\}=\mu_{p^n}.\erratum{原书误作 $(\alpha-1)^{p^n}=1$, 依 Milne errata (R. Sundararajan) 改正.}
\]
$\Lambda_n$ 上的 $+_F$ 对应于 $\mu_{p^n}$ 上的乘法, 而 $\Z_p$-模结构如前所定义. 因此 $\Lambda_n\simeq\Z/p^n\Z$ (作为 $\Z_p$-模).
\end{example}

回忆一个循环模是指可由单个元素生成的模. 由于 $A$ 是主理想整环, 且除共轭外只有一个素元, 每个有限生成挠 $A$-模 $M$ 分解为循环模的直和
\[
M\simeq A/(\pi^{n_1})\oplus\cdots\oplus A/(\pi^{n_r}),\quad n_1\le n_2\le\cdots,
\]
且序列 $n_1,\ldots,n_r$ 唯一确定.

\begin{lemma}\label{lem:3.3}
设 $M$ 是一个 $A$-模, 并令 $M_n=\Ker(\pi^n\colon M\to M)$. 假设

(a) $M_1$ 有 $q\stackrel{\mathrm{def}}{=}(A:(\pi))$ 个元素, 且

(b) $\pi\colon M\to M$ 是满的.

则 $M_n\simeq A/(\pi^n)$; 特别地, 它有 $q^n$ 个元素.
\end{lemma}

\begin{proof}
我们对 $n$ 用归纳法. 由于 $A/\pi$ 的阶为 $q$, 条件 (a) 与结构定理推出 $M_1\simeq A/(\pi)$. 考虑正合列
\[
0\to M_1\to M_n\xrightarrow{\pi}M_{n-1}\to 0.
\]
条件 (b) 表明它在 $M_{n-1}$ 处正合, 因此整体正合. 由此推出 $M_n$ 有 $q^n$ 个元素. 若 $M_n$ 不是循环的,
\[
M_n\simeq A/(\pi^{n_1})\oplus A/(\pi^{n_2})\oplus\cdots,\quad n_1,n_2,\ldots\ge 1,
\]
则 $M_1$ 也不是循环的,
\[
M_1\simeq(\pi^{n_1-1})/(\pi^{n_1})\oplus(\pi^{n_2-1})/(\pi^{n_2})\oplus\cdots.
\]
因此 $M_n$ 是一个阶为 $q^n$ 的循环 $A$-模, 而每个这样的模都同构于 $A/(\pi^n)$.
\end{proof}

\begin{proposition}\label{prop:3.4}
$A$-模 $\Lambda_n$ 同构于 $A/(\pi^n)$. 因此 $\End_A(\Lambda_n)\simeq A/(\pi^n)$, 且 $\Aut_A(\Lambda_n)\simeq(A/(\pi^n))^{\times}$.
\end{proposition}

\begin{proof}
形式群律的一个 $A$-同构 $h\colon F_f\to F_g$ 诱导 $A$-模的同构 $\Lambda_f\to\Lambda_g$, 因此我们选哪个 $f\in\mathcal{F}_{\pi}$ 无关紧要. 我们取 $f\in\mathcal{F}_{\pi}$ 为形如 $\pi T+\cdots+T^q$ 的多项式. 这是一个 Eisenstein 多项式 (ANT, p.\,129), 因此有 $q$ 个不同的根, 每个绝对值 $<1$. 设 $\alpha\in K^{\mathrm{al}}$ 绝对值 $<1$. 由
\[
f(T)-\alpha=T^q+\cdots+\pi T-\alpha
\]
的 Newton 多边形我们看到它的根绝对值 $<1$, 从而落在 $\Lambda_f$ 中. 我们已证明 $\Lambda_f$ 满足引理的假设, 故 $\Lambda_n\simeq A/(\pi^n)$. 由此推出 $A$ 在 $\Lambda_n$ 上的作用诱导一个同构 $A/(\pi^n)\to\End_A(\Lambda_n)$.
\end{proof}

\begin{lemma}\label{lem:3.5}
设 $L$ 是局部域 $K$ 的有限 Galois 扩张, Galois 群为 $G$. 对每个 $F\in\mathcal{O}_K[[X_1,\ldots,X_n]]$ 与 $\alpha_1,\ldots,\alpha_n\in\mathfrak{m}_L$,
\[
F(\tau\alpha_1,\ldots,\tau\alpha_n)=\tau F(\alpha_1,\ldots,\alpha_n),\quad\text{所有 }\tau\in G.
\]
\end{lemma}

\begin{proof}
若 $F$ 是多项式, 这由 $\tau$ 是固定 $\mathcal{O}_K$ 中元素的域同构这一事实推出. 由于 $L$ 上的绝对值是 $K$ 上绝对值的唯一扩张 (ANT, 7.38), 它被 $\tau\in G$ 保持, 即对所有 $\alpha\in L$ 有 $|\tau\alpha|=|\alpha|$, 因此 $\tau$ 连续. 于是它保持极限:
\[
\lim_{m\to\infty}\alpha_m=L\implies\lim_{m\to\infty}\tau\alpha_m=\tau L.
\]
写 $F=F_m+\text{次数}\ge m+1\text{ 的项}$, 其中 $F_m$ 是 $m$ 次多项式. 则
\[
\begin{aligned}
\tau(F(\alpha_1,\ldots,\alpha_n))&=\tau\Bigl(\lim_{m\to\infty}F_m(\alpha_1,\ldots,\alpha_n)\Bigr)\\
&=\lim_{m\to\infty}\tau\bigl(F_m(\alpha_1,\ldots,\alpha_n)\bigr)\\
&=\lim_{m\to\infty}F_m(\tau\alpha_1,\ldots,\tau\alpha_n).
\end{aligned}
\]
\end{proof}

\begin{theorem}\label{thm:3.6}
设 $K_{\pi,n}$ 表示由 $\Lambda_n$ 中元素在 $K$ 上生成的子域 $K[\Lambda_n]\subset K^{\mathrm{al}}$.

(a) 对每个 $n$, $K_{\pi,n}/K$ 是全分歧的, 次数为 $(q-1)q^{n-1}$.

(b) $A$ 在 $\Lambda_n$ 上的作用定义一个同构
\[
(A/\mathfrak{m}^n)^{\times}\to\Gal(K_{\pi,n}/K).
\]
特别地, $K_{\pi,n}/K$ 是 Abel 的.

(c) 对每个 $n$, $\pi$ 是 $K_{\pi,n}$ 中的范数.
\end{theorem}

\begin{proof}
我们再次可假定 $f\in\mathcal{F}_{\pi}$ 是形如 $\pi T+\cdots+T^q$ 的多项式. (a), (b). 选取 $f(T)$ 的一个非零根 $\pi_1$, 并 (归纳地) 选取 $f(T)-\pi_{n-1}$ 的一个根 $\pi_n$. 考虑域的序列
\[
K[\Lambda_n]\supset K[\pi_n]\stackrel{q}{\supset}K[\pi_{n-1}]\stackrel{q}{\supset}\cdots\stackrel{q}{\supset}K[\pi_1]\stackrel{q-1}{\supset}K.
\]
每个扩张都是 Eisenstein 的 (ANT, 7.55), 其次数如所标出. 因此 $K[\pi_n]$ 在 $K$ 上全分歧, 次数为 $q^{n-1}(q-1)$.

回忆 $\Lambda_n$ 是 $f^{(n)}$ 在 $K^{\mathrm{al}}$ 中的根的集合, 故 $K[\Lambda_n]$ 是 $f^{(n)}$ 的分裂域. 因此 $\Gal(K[\Lambda_n]/K)$ 可等同于 $\Lambda_n$ 上置换群的一个子群; 但引理\ref{lem:3.5} 表明 $\Gal(K[\Lambda_n]/K)$ 的每个元素在 $\Lambda_n$ 上作为一个 $A$-模同构作用, 因此 $\Gal(K[\Lambda_n]/K)$ 在 $\Lambda_n$ 的自同构群中的像包含于
\[
\End_A(\Lambda_n)=(A/(\pi^n))^{\times}.
\]
于是
\[
(q-1)q^{n-1}\ge|\Gal(K[\Lambda_n]/K)|=[K[\Lambda_n]:K]\ge[K[\pi_n]:K]=(q-1)q^{n-1}.
\]
我们必须处处取等号, 故 $\Gal(K[\Lambda_n]/K)\simeq(A/\mathfrak{m}^n)^{\times}$ 且 $K[\Lambda_n]=K[\pi_n]$.

(c) 令 $f^{[n]}(T)=(f/T)\circ f\circ\cdots\circ f$ ($n$ 项), 于是
\[
f^{[n]}(T)=\pi+\cdots+T^{(q-1)q^{n-1}}.
\]
则 $f^{[n]}(\pi_n)=f^{[n-1]}(\pi_{n-1})=\cdots=f(\pi_1)=0$. 由于 $f^{[n]}$ 是次数为 $(q-1)q^{n-1}=[K[\pi_n]:K]$ 的首一多项式, 它必是 $\pi_n$ 在 $K$ 上的极小多项式. 因此
\[
\Nm_{K[\Lambda_n]/K}\pi_n=(-1)^{(q-1)q^{n-1}}\pi=\pi\quad\text{除非 }q\text{ 为偶数且 }n=1.
\]
在例外情形 ($q$ 为偶数且 $n=1$), $\pi$ 是元素 $-\pi_1$ 的范数.\erratum{原书误作 "除非 $q=2$ 且 $n=1$", 并以 "$K[\Lambda_1]=K$" 处理例外情形; 但只要 $q$ 为偶数 (即 $2$ 的幂) 且 $n=1$, 指数 $(q-1)q^{n-1}$ 就是奇数, 此时 $\Nm_{K[\Lambda_1]/K}(\pi_1)=-\pi$, 应改取 $-\pi_1$ 的范数. 事实上对任何 $q$ 与 $n$ 都有 $\Nm_{K[\Lambda_n]/K}(-\pi_n)=\pi$, 从而无须单列例外情形. 依 Milne errata (R. Sundararajan) 改正.}
\end{proof}

\begin{summary}\label{note:3.7}
设 $f(T)=\pi T+\cdots+T^q$, 并令 $\Lambda_n$ 为 $f^{(n)}$ 在 $K^{\mathrm{al}}$ 中的根的集合. 定义 $K_{\pi,n}=K[\Lambda_n]$. 则
\[
\begin{aligned}
K_{\pi}&=\bigcup K_{\pi,n},\\
K_{\pi,n}&=K[\pi_n]\quad (f(\pi_n)=\pi_{n-1},\ \mathfrak{m}_{K_{\pi,n}}=(\pi_n)),\\
K_{\pi,2}&=K[\pi_2]\quad (f(\pi_2)=\pi_1,\ \mathfrak{m}_{K_{\pi,2}}=(\pi_2)),\\
K_{\pi,1}&=K[\pi_1]\quad (f(\pi_1)=0,\ \mathfrak{m}_{K_{\pi,1}}=(\pi_1),\ \pi_1\neq 0,\ [K_{\pi,1}:K]=q-1).
\end{aligned}
\]
此外, 作用
\[
a*\lambda=[a]_f(\lambda),\quad a\in A,\ \lambda\in\Lambda_n,
\]
诱导一个同构
\[
(A/\mathfrak{m}^n)^{\times}\to\Gal(K_{\pi,n}/K).
\]
取逆向极限, 我们得到一个同构
\[
A^{\times}\to\Gal(K_{\pi}/K).
\]
\end{summary}

\begin{example}\label{ex:3.8}
设 $K=\Q_p$, $f=(T+1)^p-1$. 对每个 $r$, 选取一个 $p^r$ 次单位根 $\zeta_{p^r}$ 使得 $\zeta_p$ 本原且 $\zeta_{p^r}^p=\zeta_{p^{r-1}}$. 则 $\pi_r=\zeta_{p^r}-1$, 且 $(\Q_p)_{p,n}=\Q_p[\pi_r]=\Q_p[\zeta_{p^r}]$. 此外同构
\[
(\Z_p/(p^n))^{\times}\to\Gal(\Q_p[\zeta_{p^r}]/\Q_p)
\]
是标准的那个.
\end{example}

\subsection*{局部 Artin 映射}

我们如下定义一个同态
\[
\phi_{\pi}\colon K^{\times}\to\Gal((K_{\pi}\cdot K^{\mathrm{un}})/K).
\]
设 $a\in K^{\times}$. 由于 $K_{\pi}\cap K^{\mathrm{un}}=K$, 只需分别描述 $\phi_{\pi}(a)$ 在 $K_{\pi}$ 与 $K^{\mathrm{un}}$ 上的作用. 写 $a=u\pi^m$, $u\in U$. 我们规定 $\phi_{\pi}(a)$ 在 $K^{\mathrm{un}}$ 上作用为 $\Frob^m$, 而它在 $K_{\pi}$ 上按如下规则作用:
\[
\phi_{\pi}(a)(\lambda)=[u^{-1}]_f(\lambda),\quad\text{所有 }\lambda\in\bigcup\Lambda_n.
\]
插入 ${}^{-1}$ 是为了使下一个定理成立.

\begin{theorem}\label{thm:3.9}
$K_{\pi}\cdot K^{\mathrm{un}}$ 与 $\phi_{\pi}$ 都与 $\pi$ 的选取无关.
\end{theorem}

本节的余下部分将用于证明此定理.

回忆 $K^{\mathrm{al}}$ 不完备 (见 ANT, 习题 7-7); 事实上即使 $K^{\mathrm{un}}$ 也不完备. 我们以 $\hat{K}^{\mathrm{un}}$ 记其完备化, 以 $B$ 记 $\hat{K}^{\mathrm{un}}$ 的赋值环 (绝对值 $|\cdot|$ 唯一地扩张到 $\hat{K}^{\mathrm{un}}$, 而 $B$ 是绝对值 $\le 1$ 的元素组成的集合). 我们以 $\sigma$ 记 $K^{\mathrm{un}}/K$ 的 Frobenius 自同构 $\Frob_K$, 也记它在 $\hat{K}^{\mathrm{un}}$ 上的扩张. 对幂级数 $\theta(T)=\sum b_i T^i\in B[[T]]$, 我们定义 $(\sigma\theta)(T)$ 为幂级数 $\sum(\sigma b_i)T^i$.

\begin{proposition}\label{prop:3.10}
设 $F_f$ 与 $F_g$ 是由 $f\in\mathcal{F}_{\pi}$ 与 $g\in\mathcal{F}_{\varpi}$ 定义的形式群律, 其中 $\pi$ 与 $\varpi=u\pi$ 是 $K$ 的两个素元. 则 $F_f$ 与 $F_g$ 在 $B$ 上成为 $A$-同构. 更精确地, 存在 $\varepsilon\in B^{\times}$ 使得 $\sigma\varepsilon=\varepsilon u$, 以及幂级数 $\theta(T)\in B[[T]]$ 使得:

(a) $\theta(T)=\varepsilon T+\text{次数}\ge 2\text{ 的项}$;

(b) $\sigma\theta=\theta\circ[u]_f$;

(c) $\theta(F_f(X,Y))=F_g(\theta(X),\theta(Y))$;

(d) $\theta\circ[a]_f=[a]_g\circ\theta$.

最后两个条件说 $\theta$ 是一个与 $A$ 的作用交换的同态 $F_f\to F_g$, 而第一个条件表明 $\theta$ 是同构 (因为 $\varepsilon$ 是单位).
\end{proposition}

\begin{lemma}\label{lem:3.11}
同态
\[
b\mapsto\sigma b-b\colon B\to B,
\]
\[
b\mapsto\sigma b/b\colon B^{\times}\to B^{\times},
\]
都是满的, 核分别为 $A$ 与 $A^{\times}$.
\end{lemma}

\begin{proof}
设 $R$ 是 $K^{\mathrm{un}}$ 中的赋值环, $\mathfrak{n}$ 是其极大理想. 则 $R$ 是离散赋值环, 且 $\varprojlim R/\mathfrak{n}^n=B$ (见下面的 (\ref{ex:A.7})). 我们对 $n$ 用归纳法证明序列
\begin{equation*}
0\longrightarrow A/\mathfrak{m}_K^n\longrightarrow R/\mathfrak{n}^n\xrightarrow{\sigma-1}R/\mathfrak{n}^n\longrightarrow 0
\tag{10}
\end{equation*}
正合. 对 $n=1$, 序列成为
\[
0\xrightarrow{\quad}k\xrightarrow{\quad}\bar{k}\xrightarrow{x\mapsto x^q-x}\bar{k}\xrightarrow{\quad}0.
\]
这里 $\bar{k}$ 是 $k$ 的代数闭包. 这显然正合. 假设序列对 $n-1$ 正合, 并考虑图表
\[
\begin{tikzcd}
0\ar[r] & R/\mathfrak{n}\ar[r]\ar[d,"\sigma-1"] & R/\mathfrak{n}^n\ar[r]\ar[d,"\sigma-1"] & R/\mathfrak{n}^{n-1}\ar[r]\ar[d,"\sigma-1"] & 0\\
0\ar[r] & R/\mathfrak{n}\ar[r] & R/\mathfrak{n}^n\ar[r] & R/\mathfrak{n}^{n-1}\ar[r] & 0.
\end{tikzcd}
\]
由蛇形引理我们发现 $\sigma-1\colon R/\mathfrak{n}^n\to R/\mathfrak{n}^n$ 是满的, 且其核有 $q^n$ 个元素. 由于 $A/\mathfrak{m}_K^n$ 包含在核中且有 $q^n$ 个元素, 它必等于该核. 这证明了 (10) 正合; 当我们对 $n$ 取逆向极限时, 序列成为
\[
0\longrightarrow A\longrightarrow B\xrightarrow{\sigma-1}B\longrightarrow 0,
\]
因而正合 (见下面的命题 \ref{prop:A.8}).

对 $B^{\times}$ 的证明类似.
\end{proof}

一个幂级数 $h$ 关于复合的逆记为 $h^{-1}$. 于是 $h\circ h^{-1}=T=h^{-1}\circ h$.

命题\ref{prop:3.10} 的证明分四步.

第 1 步. 证明存在 $\theta(T)\in B[[T]]$ 满足 (a) 与 (b).

第 2 步. 证明第 1 步中的 $\theta$ 可如此选取, 使得 $g=\sigma\theta\circ f\circ\theta^{-1}$.

第 3 步. 证明幂级数 $\theta(F_f(\theta^{-1}(X),\theta^{-1}(Y)))$ 具有刻画 $F_g(X,Y)$ 的性质, 因此等于它.\erratum{原书误作 "刻画 $F_f(X,Y)$", 依 Milne errata (A. Jangle) 改正.}

第 4 步. 证明幂级数 $\theta\circ[a]_f\circ\theta^{-1}$ 具有刻画 $[a]_g$ 的性质, 因此等于它.

\begin{proof}[第 1 步的证明]
选取 $\varepsilon\in B^{\times}$ 使得 $\sigma\varepsilon=\varepsilon u$ — 它的存在性由引理\ref{lem:3.11} 保证. 从 $\theta_1(T)=\varepsilon T$ 出发, 我们将构造一列多项式 $\theta_r$ 使得
\[
\begin{aligned}
\theta_r(T)&=\theta_{r-1}(T)+bT^r,\ \text{某个 }b\in B,\\
\sigma\theta_r&=\theta_r\circ[u]_f+\text{次数}\ge r+1\text{ 的项}.
\end{aligned}
\]
注意, 对 $\theta_1(T)=\varepsilon T$, 第二个方程成为
\[
\sigma\varepsilon T=\varepsilon(uT+\cdots)+\text{次数}\ge 2\text{ 的项},
\]
这由我们对 $\varepsilon$ 的选取为真. 设已找到 $\theta_r$, 我们要寻找 $\theta_{r+1}(T)=\theta_r(T)+bT^{r+1}$. 写 $b=a\varepsilon^{r+1}$, $a\in B$. 则第二个方程成为
\[
(\sigma\theta_r)(T)+(\sigma a)(\sigma\varepsilon)^{r+1}T^{r+1}\stackrel{?}{=}\theta_r([u]_f(T))+a\varepsilon^{r+1}(uT)^{r+1}+\text{次数}\ge r+1\text{ 的项}.
\]
因此我们需要
\[
(\sigma a-a)(\varepsilon u)^{r+1}\stackrel{?}{=}c,
\]
其中 $c$ 是 $\theta_r\circ[u]_f-\sigma\theta_r$ 中 $T^{r+1}$ 的系数. 我们可以取 $a$ 为 $B$ 中任何满足 $\sigma a-a=c/(\varepsilon u)^{r+1}$ 的元素.
\end{proof}

\begin{proof}[第 2 步的证明]
定义
\[
h=\sigma\theta\circ f\circ\theta^{-1}=\theta\circ[u]_f\circ f\circ\theta^{-1}=\theta\circ f\circ[u]_f\circ\theta^{-1}.
\]
则由于 $f$ 与 $[u]_f$ 的系数在 $A$ 中,
\[
\sigma h=\sigma\theta\circ f\circ[u]_f\circ\sigma\theta^{-1}=\sigma\theta\circ f\circ\theta^{-1}=h.
\]
中间等号用了 $[u]_f\circ\sigma\theta^{-1}=\theta^{-1}$, 这由 $\theta\circ[u]_f\circ\sigma\theta^{-1}=T$ 推出. 由于 $\sigma h=h$, 它落在 $A[[T]]$ 中. 此外,
\[
h(T)=\sigma\varepsilon\cdot\pi\cdot\varepsilon^{-1}T+\cdots=\varpi T+\text{次数}\ge 2\text{ 的项},
\]
且
\[
h(T)\equiv\sigma\theta\circ(\theta^{-1})^q\equiv\sigma\theta(\sigma\theta^{-1}(T^q))\equiv T^q\bmod\mathfrak{m}_K.
\]
因此 $h\in\mathcal{F}_{\varpi}$. 令 $\theta'=[1]_{g,h}\circ\theta$. 则 $\theta'$ 显然仍满足命题的条件 (a), 且仍满足 (b), 因为 $[1]_{g,h}\in A[[T]]$. 此外,
\[
\sigma\theta'\circ f\circ\theta'^{-1}=[1]_{g,h}\circ h\circ[1]_{g,h}^{-1}=g.
\]
\end{proof}

第 3 步与第 4 步的证明是引理\ref{lem:2.11} 的直接应用, 留给读者.

\begin{proof}[$K_{\pi}\cdot K^{\mathrm{un}}$ 与 $\pi$ 无关的证明]
设 $\pi$ 与 $\varpi=\pi u$ 是 $K$ 的两个素元. 由命题\ref{prop:3.10} 我们得到
\[
(\sigma\theta)\circ[\pi]_f=\theta\circ[u]_f\circ[\pi]_f=\theta\circ[\varpi]_f=[\varpi]_g\circ\theta,
\]
即
\[
(\sigma\theta)(f(T))=g(\theta(T)).
\]
因此, 对任何 $\alpha\in K^{\mathrm{al}}$ (回忆这是 $K$ 的可分代数闭包),
\[
f(\alpha)=0\implies g(\theta(\alpha))=0,
\]
且类似地
\[
g(\alpha)=0\implies f(\theta^{-1}(\alpha))=0.
\]
因此 $\theta$ 定义双射 $\Lambda_{f,1}\to\Lambda_{g,1}$, 从而
\[
\hat{K}^{\mathrm{un}}[\Lambda_{g,1}]=\hat{K}^{\mathrm{un}}[\theta(\Lambda_{f,1})]\subset\hat{K}^{\mathrm{un}}[\Lambda_{f,1}]=\hat{K}^{\mathrm{un}}[\theta^{-1}(\Lambda_{g,1})]\subset\hat{K}^{\mathrm{un}}[\Lambda_{g,1}].
\]
因此
\[
\hat{K}^{\mathrm{un}}[\Lambda_{g,1}]=\hat{K}^{\mathrm{un}}[\Lambda_{f,1}].
\]
现在下一个引理表明
\[
\hat{K}^{\mathrm{un}}[\Lambda_{g,1}]\cap K^{\mathrm{al}}=K^{\mathrm{un}}[\Lambda_{g,1}],\qquad \hat{K}^{\mathrm{un}}[\Lambda_{f,1}]\cap K^{\mathrm{al}}=K^{\mathrm{un}}[\Lambda_{f,1}],
\]
于是
\[
K^{\mathrm{un}}[\Lambda_{g,1}]=K^{\mathrm{un}}[\Lambda_{f,1}].
\]
论证可无困难地推广, 表明对所有 $n$ 有
\[
K^{\mathrm{un}}[\Lambda_{g,n}]=K^{\mathrm{un}}[\Lambda_{f,n}],
\]
从而 $K^{\mathrm{un}}\cdot K_{\varpi}=K^{\mathrm{un}}\cdot K_{\pi}$.
\end{proof}

\begin{lemma}\label{lem:3.12}
$K^{\mathrm{al}}$ 中每个含 $K$ 的子域 $E$ 都是闭的 (在拓扑意义下).
\end{lemma}

\begin{proof}
令 $H=\Gal(K^{\mathrm{al}}/E)$. 则 $H$ 固定 $E$ 的每个元素, 因此由连续性它也固定 $E$ 的闭包中的每个元素. 按 Galois 理论, 这表明 $E$ 等于它在 $K^{\mathrm{al}}$ 中的闭包.
\end{proof}

\begin{proof}[$\phi_{\pi}$ 与 $\pi$ 无关的证明]
我们将证明, 对任意两个素元 $\pi$ 与 $\varpi$,
\[
\phi_{\pi}(\varpi)=\phi_{\varpi}(\varpi).
\]
由于 $\pi$ 任意, 这推出对任何其他的素元 $\pi'$,
\[
\phi_{\pi'}(\varpi)=\phi_{\varpi}(\varpi)=\phi_{\pi}(\varpi).
\]
由于 $\varpi$ 也任意, 且素元生成群 $K^{\times}$, 这就推出 $\phi_{\pi}=\phi_{\pi'}$.

在 $K^{\mathrm{un}}$ 上, $\phi_{\pi}(\varpi)$ 与 $\phi_{\varpi}(\varpi)$ 都诱导 Frobenius 自同构. 剩下要证明它们在 $K_{\varpi}$ 上有相同的作用.

设 $\theta$ 是如命题\ref{prop:3.10} 中在 $\hat{K}^{\mathrm{un}}$ 上的同构 $F_f\to F_g$. 它对所有 $n$ 诱导同构 $\Lambda_{f,n}\to\Lambda_{g,n}$. 按定义, $\phi_{\varpi}(\varpi)$ 在 $K_{\varpi}$ 上是恒等; 由于 $K_{\varpi,n}$ 由元素 $\theta(\lambda)$ ($\lambda\in\Lambda_{f,n}$) 在 $K$ 上生成, 剩下只需证明
\[
\phi_{\pi}(\varpi)(\theta(\lambda))=\theta(\lambda),\qquad\text{所有 }\lambda\in\Lambda_{f,n}.
\]
写 $\varpi=u\pi$. 则 $\phi_{\pi}(\varpi)=\phi_{\pi}(u)\cdot\phi_{\pi}(\pi)=\tau\sigma$, 其中
\[
\sigma=\begin{cases}\Frob_K&\text{在 }K^{\mathrm{un}}\text{ 上},\\ \mathrm{id}&\text{在 }\lambda\text{ 上},\end{cases}\qquad
\tau=\begin{cases}\mathrm{id}&\text{在 }K^{\mathrm{un}}\text{ 上},\\ [u^{-1}]_f&\text{在 }\lambda\text{ 上}.\end{cases}
\]
利用级数 $\theta$ 的系数在 $\hat{K}^{\mathrm{un}}$ 中以及 (\ref{prop:3.10}), 我们发现
\[
\phi_{\pi}(\varpi)(\theta(\lambda))=\tau\sigma(\theta(\lambda))=(\sigma\theta)(\tau\lambda)=(\sigma\theta)([u^{-1}]_f(\lambda))=\theta(\lambda).
\]
\end{proof}

\begin{example}\label{ex:3.13}
我们描述局部 Artin 映射 $\phi_p\colon\Q_p^{\times}\to\Gal(L/\Q_p)$ 在 $L=\Q_p[\zeta]$ 的情形, 其中 $\zeta$ 是一个本原 $n$ 次单位根.

(a) 设 $n$ 与 $p$ 互素. 则 $L$ 在 $\Q_p$ 上非分歧, 其次数等于剩余类域扩张的次数. 剩余类域是 $\F_{p^f}$, 其中 $p^f$ 是使得 $n\mid(p^f-1)$ 的 $p$ 的最小幂. 映射 $\phi_p\colon\Q_p^{\times}\to\Gal(L/\Q_p)$ 把 $u\cdot p^t$ 送到 Frobenius 元素的 $t$ 次幂, 其核为 $\Z_p^{\times}\cdot\langle p^f\rangle$.

(b) 设 $n$ 是 $p$ 的幂 $p^r$. 此时 $L$ 在 $K$ 上全分歧, 次数为 $(p-1)p^{r-1}$, 且 $L=(\Q_p)_{p,n}$ (见 \ref{ex:3.8}). 映射 $\phi_p\colon\Q_p^{\times}\to\Gal(L/\Q_p)$ 可描述如下: 令 $a=up^t$, 并令 $u_0\in\Z$ 表示 $u$ 在 $(\Z_p/p^r\Z_p)^{\times}$ 中的类; 则 $\phi_p(a)$ 把 $\zeta$ 送到 $\zeta^{u_0^{-1}}$, 其中 $u_0$ 是一个模 $p^r$ 同余于 $u$ 的整数. 其核为 $\{up^m\mid u\equiv 1\bmod p^r,\ m\in\Z\}$.

(c) 一般情形, 写 $n=m\cdot p^r$, 其中 $m$ 与 $p$ 互素. 则 $L=\Q_p[\zeta_m]\cdot\Q_p[\zeta_{p^r}]$, 且映射
\[
\Q_p^{\times}/\Nm(\Q_p[\zeta_n]^{\times})\to\Gal(\Q_p[\zeta_n]/\Q_p)
\]
可描述如下: 写 $a=up^t$, $u\in\Z_p^{\times}$; 则 $a$ 在 $\Q_p[\zeta_m]$ 上作用为 $\zeta_m\mapsto\zeta_m^{p^t}$,\erratum{原书误作 $\zeta_m\mapsto\zeta_m^{t}$, 依 Milne errata (R. Sundararajan) 改正.} 而在 $\Q_p[\zeta_{p^r}]$ 上作用为 $\zeta_{p^r}\mapsto\zeta_{p^r}^{u_0^{-1}}$, 其中 $u_0$ 是一个模 $p^r$ 同余于 $u$ 的整数.
\end{example}

\section{局部 Kronecker-Weber 定理}\label{sec:I.4}

本节证明的主要结果是 $K^{\mathrm{ab}}=K_{\pi}\cdot K^{\mathrm{un}}$. 由于局部类域论主定理的证明并不需要它, 而且它被这些定理所蕴含, 本节可以跳过.

\subsection*{$K_{\pi,n}/K$ 的分歧群}

设 $L/K$ 是有限 Galois 扩张, Galois 群为 $G$. 回忆 (ANT, 7.57, 7.58, \ldots) 第 $i$ 个分歧群定义为
\[
G_i=\{\tau\in G\mid\mathrm{ord}_L(\tau a-a)\ge i+1,\ \text{所有 }a\in\mathcal{O}_L\};
\]
此外, 对 $i\ge 1$,
\[
G_i=\{\tau\in G_0\mid\mathrm{ord}_L(\tau\varPi-\varPi)\ge i+1\},
\]
其中 $\varPi$ 是 $L$ 的一个素元. 这里 $\mathrm{ord}_L$ 是正规化赋值 $L^{\times}\to\Z$. 于是 $G/G_0=\Gal(l/k)$, 并有包含关系:
\[
\begin{aligned}
(\varPi\mapsto\tau\varPi/\varPi\bmod\varPi)&\colon G_0/G_1\hookrightarrow l^{\times},\\
(\varPi\mapsto(\tau\varPi-\varPi)/\varPi^{i+1}\bmod\varPi)&\colon G_i/G_{i+1}\hookrightarrow l,
\end{aligned}
\]
其中 $l\supset k$ 是 $L$ 与 $K$ 的剩余类域. 因此 $(G_0:G_1)\mid \#l-1$, 且对 $i\ge 1$ 有 $(G_i:G_{i+1})\mid \#l$ (当 $L/K$ 全分歧时 $l=k$, 从而 $\#l=q$).\erratum{原书此处直接写 $q$; 依 Milne errata (Zhang, Zhuo), 应改为 $l$ 的阶, 它只在引理\ref{lem:4.9} 的情形 (全分歧) 才等于 $q=|k|$.} 此外对充分大的 $i$ 有 $G_i=\{1\}$.

令
\[
U^{(0)}=U=A^{\times},\qquad U^{(i)}=1+\mathfrak{m}^i,\ i\ge 1.
\]
则在 $A^{\times}/(1+\mathfrak{m}^n)=U/U^{(n)}$ 上我们有一个滤列
\[
U/U^{(n)}\supset U^{(1)}/U^{(n)}\supset\cdots\supset U^{(n)}/U^{(n)}=0.
\]

\begin{proposition}\label{prop:4.1}
设 $L=K_{\pi,n}$. 在定理\ref{thm:3.6} 的同构 $A^{\times}/U^{(n)}\xrightarrow{\simeq}G$ 下, $U^{(i)}/U^{(n)}$ 映满到 $G_{q^i-1}$.
\end{proposition}

\begin{proof}
我们取 $f=\pi T+T^q$. 显然 $G=G_0$, 且 $U^{(0)}/U^{(n)}$ 映到 $G_0$. 现在取 $i\ge 1$, 并设 $u\in U^{(i)}\setminus U^{(i+1)}$. 则 $u=1+v\pi^i$, 其中 $v\in A^{\times}$, 且
\[
[u]_f(\pi_n)=[1]_f(\pi_n)+[v]_f[\pi^i]_f(\pi_n)=\pi_n+[v]_f(\pi_{n-i})=\pi_n+(\text{单位})\pi_{n-i}.
\]
对每个 $i\ge 1$, $\pi_i=\pi\pi_{i+1}+\pi_i^q=\pi_{i+1}^q(\frac{\pi\pi_{i+1}}{\pi_{i+1}^q}+1)=\pi_{i+1}^q\times\text{单位}$, 因为 $\mathrm{ord}(\frac{\pi}{\pi_{i+1}^{q-1}})>0$. 因此 $\pi_{n-i}=\pi_n^{q^i}\times\text{单位}$, 且
\[
[u]_f(\pi_n)-\pi_n=\pi_n^{q^i}\times\text{单位}.
\]
按定义, 这意味着 $[u]_f\in G_{q^i-1}$, $[u]_f\notin G_{q^i}$. 由于这对所有 $i$ 成立, 它推出 $U^{(i)}$ 映到 $G_{q^i-1}$.

于是
\[
\begin{aligned}
G_0&=G,\\
G_{q-1}&=G_{q-2}=\cdots=G_1,\\
G_{q^2-1}&=G_{q^2-2}=\cdots=G_q,\\
&\ \vdots\\
G_{q^n-1}&=1
\end{aligned}
\]
是 $K_{\pi,n}/K$ 的一组完整的互异分歧群.
\end{proof}

\subsection*{分歧群的上标号}

设 $L$ 是 $K$ 的有限 Galois 扩张, Galois 群为 $G$. 我们把 $G_u$ 的定义扩充到所有实数 $u\ge -1$, 令
\[
G_u=G_i,\quad\text{其中 }i\text{ 是 }\ge u\text{ 的最小整数}.
\]
对 $u>0$,
\[
G_u=\{\tau\in G_0\mid\mathrm{ord}_L(\tau\varPi-\varPi)\ge i+1\}.
\]
定义 $\varphi\colon\R_{\ge 0}\to\R$ 为唯一的连续分段线性函数使得
\[
\begin{cases}
\varphi(0)=0,\\
\varphi'(u)=(G_0:G_u)^{-1}\quad\text{若 }u\text{ 不是整数}.
\end{cases}
\]
定义 $G^v=G_u$, 若 $v=\varphi(u)$, 即 $G^v=G_{\varphi^{-1}(v)}$.

\begin{example}\label{ex:4.2}
设 $L=K_{\pi,n}$. 则
\[
(G_0:G_1)=q-1,\qquad G_1=G_2=\cdots=G_{q-1}.
\]
因此 $\varphi'(u)=\frac{1}{q-1}$ 对 $0<u<q-1$, 且 $\varphi$ 的图像第一段从 $(0,0)$ 跑到 $(q-1,1)$; 故 $G^1=G_{q-1}$. 接着
\[
(G_{q-1}:G_q)=q,\qquad G_q=G_{q+1}=\cdots=G_{q^2-1}.
\]
因此 $\varphi'(u)=\frac{1}{q(q-1)}$ 对 $q-1<u<q^2-1$, 且 $\varphi$ 的图像第二段从 $(q-1,1)$ 跑到 $(q^2-1,2)$. 于是 $G^2=G_{q^2-1}$. 继续这样做, 我们得到如下图像:
\[
\begin{tikzcd}[column sep=small, row sep=small]
G_0 \ar[d, no head, double] & \overset{q-1}{\supset} & G_{q-1} \ar[d, no head, double] & \overset{q}{\supset} & G_{q^2-1} \ar[d, no head, double] & \supset & \cdots & \supset & G_{q^n-1}=1 \ar[d, no head, double] \\
G^0 & & G^1 & & G^2 & & \cdots & & G^n
\end{tikzcd}
\]
\end{example}

\begin{proposition}\label{prop:4.3}
在同构 $A^{\times}/U^{(n)}\to G$ 下,
\[
U^{(i)}/U^{(n)}\xrightarrow{\approx}G^i.
\]
\end{proposition}

\begin{proof}
由 (\ref{prop:4.1}) 与 (\ref{ex:4.2}) 直接推出.
\end{proof}

上标号的定义使得它与过渡到商相容 (而下标号与过渡到子群相容).

\begin{proposition}\label{prop:4.4}
考虑 Galois 扩张 $M\supset L\supset K$, 并令 $G=\Gal(M/K)$, $H=\Gal(M/L)$ (假定为正规的), 于是 $G/H=\Gal(L/K)$. 则
\[
(G/H)^v=\mathrm{Im}(G^v\to G/H),
\]
即 $(G/H)^v=G^vH/H$.
\end{proposition}

\begin{proof}
见 \cite{Serre1962}, IV.3, 命题 14.
\end{proof}

现在考虑一个无限 Galois 扩张 $\Omega/K$. 利用 (\ref{prop:4.4}) 我们可以在 $G=\Gal(\Omega/K)$ 上定义一个滤列:
\[
\tau\in G^v\iff\tau\in\Gal(L/K)^v,\quad\text{所有有限 Galois 子扩张 }L\subset\Omega.
\]

\begin{definition}\label{def:4.5}
对有限 Galois 扩张 $L/K$, 若对所有 $\varepsilon>0$ 都有 $G^v\neq G^{v+\varepsilon}$, 则称 $v$ 是滤列 $\{G^v\}$ 中的一个\emph{跳跃}\index{fenqiqun@分歧群(ramification group)!tiaoyue@跳跃(jump)}.
\end{definition}

\begin{theorem}[Hasse-Arf]\label{thm:4.6}
若 $L/K$ 是 Abel 的, 则跳跃都是整数, 即若 $G_i\neq G_{i+1}$, 则 $\varphi(i)\in\Z$.
\end{theorem}

\begin{proof}
见 \cite{Serre1962}, V.7. (证明相当初等但复杂. 它不要求剩余类域有限, 只要求剩余类域扩张可分.)
\end{proof}

因此, 对有限 Abel 扩张 $L/K$, $G_0=G^0$ 上的滤列形如
\[
G^0\supsetneq G^{i_1}\supsetneq G^{i_2}\ldots\quad i_j\in\N.
\]
此外对充分大的 $n$ 有 $G^n=\{1\}$, 且 $(G^{i_j}:G^{i_j+1})$ 整除 $q-1$ 或 $q$. 对无限 Abel 扩张, 同样的陈述成立, 只是滤列不必终止: 我们只能说
\[
\bigcap G^i=\{1\}.
\]

\begin{example}\label{ex:4.7}
设 $L=K_{\pi,n}$. 若 $G_i\neq G_{i+1}$, 则 $i=0,q-1,\ldots,q^n-1$, 而在那些点 $\varphi$ 取值 $1,2,3,\ldots,n$. 因此我们对所有这些扩张验证了 Hasse-Arf 定理, 并且由 (\ref{prop:4.4}), 对所有子扩张也成立. 特别地, 我们表明了局部 Kronecker-Weber 定理蕴含 Hasse-Arf 定理.
\end{example}

\subsection*{局部 Kronecker-Weber 定理}

如常, $K$ 是局部非阿基米德 域, 且 $K$ 的所有扩张都要求为某个固定的可分代数闭包 $K^{\mathrm{al}}$ 的子域.

\begin{theorem}\label{thm:4.8}
对 $K$ 的每个素元 $\pi$,
\[
K_{\pi}\cdot K^{\mathrm{un}}=K^{\mathrm{ab}}.
\]
\end{theorem}

\begin{lemma}\label{lem:4.9}
设 $L$ 是 $K$ 的 Abel 全分歧扩张. 若 $L\supset K_{\pi}$, 则 $L=K_{\pi}$.
\end{lemma}

\begin{proof}
令 $G=\Gal(L/K)$, $H=\Gal(L/K_{\pi})$, 于是 $G/H=\Gal(K_{\pi}/K)$. 考虑 (Abel 群的) 图表
\[
\begin{tikzcd}[column sep=small]
 & 1\ar[d] & 1\ar[d] & 1\ar[d] & \\
1\ar[r] & G^{n+1}\cap H\ar[r]\ar[d] & G^{n+1}\ar[r]\ar[d] & (G/H)^{n+1}\ar[r]\ar[d] & 1\\
1\ar[r] & G^n\cap H\ar[r]\ar[d] & G^n\ar[r]\ar[d] & (G/H)^n\ar[r]\ar[d] & 1\\
1\ar[r] & \frac{G^n\cap H}{G^{n+1}\cap H}\ar[r]\ar[d] & \frac{G^n}{G^{n+1}}\ar[r]\ar[d] & \frac{(G/H)^n}{(G/H)^{n+1}}\ar[r]\ar[d] & 1\\
 & 1 & 1 & 1 &
\end{tikzcd}
\]
其中各列显然正合, 而命题\ref{prop:4.4} 表明上面两行正合. 因此 (例如由蛇形引理) 第三行也正合, 从而
\[
(G^n:G^{n+1})=((G/H)^n:(G/H)^{n+1})(G^n\cap H:G^{n+1}\cap H).
\]
由于右边 $\le q$ 且 $(G/H)^n:(G/H)^{n+1}$ 等于 $q-1$ 或 $q$, 我们推出对所有 $n$ 有 $G^n\cap H=G^{n+1}\cap H$. 于是
\[
G^{n+1}\cap H=G^n\cap H=\cdots=G^0\cap H=H,
\]
即对所有 $n$ 有 $H\subset G^{n+1}$. 由于 $\bigcap G^n=1$, 这表明 $H=1$.
\end{proof}

\begin{lemma}\label{lem:4.10}
$K_{\pi}$ 的每个有限非分歧扩张都包含在 $K_{\pi}\cdot K^{\mathrm{un}}$ 中.
\end{lemma}

\begin{proof}
设 $L$ 是 $K_{\pi}$ 的一个非分歧扩张. 则对某个 $n$, $L=K_{\pi}\cdot L'$ 对某 $K_{\pi,n}$ 的非分歧扩张 $L'$ 成立.\footnote{$K$ 的有限扩张 $L$ 是可分的, 故由某个元素 $\theta$ 生成, 而 $\theta$ 在 $K$ 上的极小方程只有有限多个系数. 由于 $K_{\pi}$ 是子域链 $K_{\pi,n}$ 的并, 其中必有一个含有该多项式的全部系数。} 现在应用 (ANT, 7.58) 可知 $L'=K_{\pi,n}\cdot L''$ 对某 $K$ 的非分歧扩张 $L''$ 成立.
\end{proof}

\begin{lemma}\label{lem:4.11}
设 $L$ 是 $K_{\pi}$ 的一个指数为 $m$ 的有限 Abel 扩张 (即对所有 $\tau\in\Gal(L/K)$, $\tau^m=1$), 并设 $K_m$ 是 $K_{\pi}$ 的 $m$ 次非分歧扩张. 则存在 $K_{\pi}$ 的一个全分歧 Abel 扩张 $L_t$ 使得
\[
L\subset L_t\cdot K_m=L\cdot K_m.
\]
\end{lemma}

\begin{proof}
对每个 $\tau\in\Gal(LK_m/K_{\pi})$, $\tau^m|_L=1=\tau^m|_{K_m}$, 因此 $\Gal(LK_m/K_{\pi})$ 仍是指数为 $m$ 的 Abel 群. 取 $\tau\in\Gal(LK_m/K_{\pi})$ 使得 $\tau|_{K_m}$ 是 Frobenius 自同构. 则 $\tau$ 的阶为 $m$, 因此
\[
\Gal(LK_m/K)=\langle\tau\rangle\times H\quad(\text{直积})
\]
对某子群 $H$. 令 $L_t=L^{\langle\tau\rangle}$; 则 $L_t$ 在 $K$ 上全分歧, 且 $L\cdot K_m=L_t\cdot K_m$.
\end{proof}

\begin{proof}[定理\ref{thm:4.8} 的证明]
设 $L$ 是 $K$ 的有限 Abel 扩张; 我们要证明 $L\subset K_{\pi}\cdot K^{\mathrm{un}}$. 把引理\ref{lem:4.11} 用于 $L\cdot K_{\pi}$ 表明存在 $K_{\pi}$ 的全分歧扩张 $L_t$ 与 $K_{\pi}$ 的非分歧扩张 $L_u$ 使得
\[
L\cdot K_{\pi}\subset L_t\cdot L_u.
\]
现在 (\ref{lem:4.9}) 推出 $L_t\subset K_{\pi}$, 而 (\ref{lem:4.10}) 推出 $L_u\subset K_{\pi}\cdot K^{\mathrm{un}}$.
\end{proof}

\begin{corollary}\label{cor:4.12}
$\Q_p$ 的每个有限 Abel 扩张都包含在一个分圆扩张中.
\end{corollary}

\begin{example}\label{ex:4.13}
$K$ 的一个有限 Abel 扩张 $L$ 未必能写成 $L=L_t\cdot L_u$, 其中 $L_t$ 在 $K$ 上全分歧, $L_u$ 在 $K$ 上非分歧. 考虑: 域 $\Q_5[\zeta_5]$ 在 $\Q_5$ 上全分歧, 次数为 $4$, Galois 群 $(\Z/5)^{\times}$ 是阶 $4$ 的循环群. 注意 $624=5^4-1$, 因此 $\Q_5[\zeta_{624}]$ 在 $\Q_5$ 上非分歧, 次数为 $4$, 其 Galois 群也是循环的. 显然
\[
\Gal(\Q_5[\zeta_{3120}]/\Q_5)=\langle\sigma\rangle\times\langle\tau\rangle,
\]
其中
\[
\begin{cases}
\sigma(\zeta_{624})=\zeta_{624}^5,\\
\sigma(\zeta_5)=\zeta_5,
\end{cases}\qquad
\begin{cases}
\tau(\zeta_{624})=\zeta_{624},\\
\tau(\zeta_5)=\zeta_5^2.
\end{cases}
\]
令 $L$ 为 $\langle\sigma^2\tau\rangle$ 的固定域. 则 $L$ 是 $\Q_5$ 的循环 Galois 扩张, 次数为 $4$. 它的最大非分歧子域
\[
L_u=L\cap\Q_5[\zeta_{624}]=\Q_5[\zeta_{624}]^{\langle\sigma^2\rangle}=\Q_5[\zeta_{624}^{25}]
\]
在 $\Q_5$ 上次数为 $2$. 若存在域 $L_t$ 使得 $L=L_t\cdot L_u$, 则 $\Gal(L/\Q_5)$ 将是两个阶 $2$ 循环群的乘积, 与它是循环群矛盾.

我们无需命题\ref{prop:3.10} 就重新得到了 $K_{\pi}\cdot K^{\mathrm{un}}$ 与 $\pi$ 无关的结论. 然而, 要证明 $\phi_{\pi}$ 与 $\pi$ 无关仍需要该命题.
\end{example}

\begin{remark}\label{rem:4.14}
原始的 Kronecker-Weber 定理 (由 Hilbert 在 1896 年利用对分歧群的分析证明, 此前 Kronecker 与 Weber 有不完全的证明) 陈述: $\Q$ 的每个有限 Abel 扩张都包含在一个分圆扩张中. 对 $\Q_p$ 同样的陈述称为\emph{局部 Kronecker-Weber 定理}, 而定理\ref{thm:4.8} 通常被称为 $K$ 上的局部 Kronecker-Weber 定理. 事实上可以对 $\Q_p$ 给出局部 Kronecker-Weber 定理的一个初等证明 (见 \cite{Cassels1986}, p.\,151).
\end{remark}

\begin{remark}\label{rem:4.15}
在第 \ref{ch:III} 章, 我们将不借助 Hasse-Arf 定理、直接从定理\ref{thm:1.1} 推出局部 Kronecker-Weber 定理 — 这是 Lubin 与 Tate 的原始思路. 上面的证明遵循 \cite{Gold1981} 与 \cite{Lubin1981} (显然 Leopoldt 更早发现了类似的证明但未发表). 对特征零的局部域, 不用 Hasse-Arf 定理或上同调、但比上述更复杂的局部 Kronecker-Weber 定理证明, 见 \cite{Rosen1981}. 正如 Iwasawa 所指出的 (1986, p.\,115), 一旦命题\ref{prop:4.4} 与 Abel 扩张 $K_{\pi,n}/K$ 的某些性质被当作已知, 那么 $K$ 上的局部 Kronecker-Weber 定理与 Hasse-Arf 定理本质上等价.
\end{remark}

\subsection*{整体 Kronecker-Weber 定理}

既然现在如此容易, 我们不妨也证明原始的 Kronecker-Weber 定理.

\begin{theorem}\label{thm:4.16}
$\Q$ 的每个 Abel 扩张都包含在一个分圆扩张中.
\end{theorem}

\begin{lemma}\label{lem:4.17}
设 $K$ 是 $\Q$ 的有限 Galois 扩张, Galois 群为 $G$. 则 $G$ 由那些在扩张 $K/\Q$ 中分歧的 $K$ 的素理想 $\mathfrak{p}$ 的惯性群生成.
\end{lemma}

\begin{proof}
令 $H$ 为惯性群生成的 $G$ 的子群, 并令 $M$ 为 $H$ 的固定域. 则对所有 $K$ 的素理想 $\mathfrak{p}$ 有 $K^{I(\mathfrak{p})}\supset M$, 因此 $\mathfrak{p}\cap M$ 在扩张 $M/\Q$ 中非分歧. 于是 $M$ 是 $\Q$ 的非分歧扩张, 从而等于 $\Q$ (由 ANT, 4.9).
\end{proof}

\begin{proof}[Kronecker-Weber 定理的证明]\footnote{这里遵循 \cite{Cassels1986}, p.~236。}
设 $K$ 是 $\Q$ 的 Abel 扩张. 设 $p$ 为一个素数, 并设 $\mathfrak{p}$ 是 $\mathcal{O}_K$ 中位于其上的素理想. 由局部 Kronecker-Weber 定理, $K_{\mathfrak{p}}$ 包含在一个 $\Q_p$ 的分圆扩张中, 记为 $K_{\mathfrak{p}}=\Q_p[u_p,v_p]$, 其中 $u_p$ 是一个 $p^{s_p}$ 次单位根, $v_p$ 是一个与 $p$ 互素次数的单位根. 注意 $s_p$ 只依赖于 $p$ (而非 $\mathfrak{p}$), 因为 $\Gal(K_\mathfrak{p}/\Q_p)$ 在位于 $p$ 之上的素理想上传递作用, 从而在域 $K_\mathfrak{p}$ 的集合上传递.

令 $L$ 为 $\Q$ 的由所有在 $K$ 中分歧的素数 $p$ 的 $p^{s_p}$ 次单位根生成的分圆扩张, 并令 $K'=K\cdot L$. 则 $K'$ 在 $\Q$ 上再次 Abel, 且对 $K'$ 的每个素理想 $\mathfrak{p}'$, 我们有
\begin{equation*}
K_{\mathfrak{p}'}'\subset\Q_p[u_p,w_p],
\tag{$*$}
\end{equation*}
其中 $w_p$ 是一个与 $p$ 互素次数的单位根. 显然用 $K'$ 替换 $K$ 来证明定理已足够, 因此我们可以假设 $K\supset L$.

现在我们有
\[
[K:\Q]\ge[L:\Q]=\prod_p\varphi(p^{s_p})\quad(\text{乘积取遍在 }K\text{ 中分歧的素数 }p).
\]
由于 $K/\Q$ 的 Galois 群 $G$ 是交换的, 惯性群 $I(p)$ 只依赖于底层的素数 $p$, 因此记为 $I(p)$. 由 (*) 我们有
\[
(I(p):1)\le\varphi(p^{s_p})
\]
因为 $I(p)$ 是 $\Q_p[u_p,w_p]$ 相应群的商群. 由 (\ref{lem:4.17}), $G$ 由群 $I(p)$ 生成, 因此存在满射 $\prod I(p)\to G$; 于是
\[
(G:1)\le\prod_p(I(p):1)\le\prod\varphi(p^{s_p}).
\]
但 $(G:1)=[K:\Q]$, 因此处处取等号, 且 $K=L$.
\end{proof}

\begin{remark}\label{rem:4.18}
在这一点上, 完成整体类域论主定理在 $K=\Q$ 情形的证明已不太困难 (见引言). 由 $L$ 包含在一个分圆扩张这一事实我们推出 Artin 映射 $\phi_L/\Q$ 有一个模. 现在用 Dirichlet 关于算术级数中素数密度的定理表明 $\phi_L/\Q$ 是满的. 我们知道 $\Nm(I_L^S)$ 包含在核中, 因此此刻所缺的只是
\[
(I_K^S:\Nm I_L^S\cdot i(K_{\mathfrak{m},1}))\le[L:\Q].
\]
这是第一个不等式, 用分析方法不难 (见 \cite{Janusz1996}, IV, 5.6). 一旦有了它, 存在定理就由 $\Q[\zeta_m]$ 有模 $(m)\infty$ 这一事实推出.
\end{remark}

\begin{NOTES}\label{note:4.19}
Lubin-Tate 扩张理论的原始文献是 Lubin 与 Tate \citeyear{LubinTate1965}. 该理论也在 \cite{Serre1967a}, \cite{Iwasawa1986}, 以及 Fesenko 与 Vostokov \citeyear{FesenkoVostokov1993} 中论述.
\end{NOTES}

\subsection*{这一切从何而来?}\label{note:4.20}
我们看到 Lubin-Tate 形式群律对一个表面上非常复杂的问题 — 给出局部域最大 Abel 扩张的显式生成元并描述 Galois 群如何在它们上作用 — 提供了一个异常简单的解. 唯一的神秘之处是有人怎么会想到它们. 下面的揣测或许有帮助.

这种理论可能存在, 是由椭圆曲线的复乘理论暗示的. 这里人们证明, 对二次虚数域 $K$, 存在一条椭圆曲线 $E$, 在同构意义下唯一, 以 $\mathcal{O}_K$ 为其自同态环. 对每个 $n$, $E$ 上阶整除 $n$ 的点构成一个循环 $\mathcal{O}_K$-模, 正是这一事实使得添加它们的坐标得到 Abel 扩张成为可能.

在为非阿基米德局部域寻找类似理论时, 自然要用局部类比 — 即形式群律 — 来替换代数群 $E$. 因此我们寻找一个形式群律, 其自同态环大到足以使它的挠点构成一个循环模. 自同态环显而易见的候选仍然是 $K$ 中的整数环 $\mathcal{O}_K$. 起初自然只要求该形式群承认一个对应于 $\mathcal{O}_K$ 的素元 $\pi$ 的自同态. 对形式群律高度 (height) 的考虑, 加上希望挠点构成循环模, 暗示这个自同态应由一个满足 $f(T)\equiv T^q\bmod\mathfrak{m}$ 的幂级数 $f(T)$ 给出. 此外, 由于我们确实希望该形式群依赖于 $\pi$ 的选取 (因为我们想构造的扩张 $K_{\pi}$ 依赖于它), 自然要求 $f(T)=\pi T+\cdots$. 于是我们被引向集合 $\mathcal{F}_{\pi}$, 而一旦有了它, 理论便自然展开.

实际上, 真实的故事更有趣 — 见下.

\begin{namedblock}{LUBIN MO220796}
由于我曾读过并欣赏 Lazard 关于一维形式群(律)的论文——该文处理的是基域特征为 $p$ 的情形——我决定考察 $p$-进环上的形式群. 出于某种原因, 我想了解这些东西的自同态环, 并逐渐认识到两类情形的相似性: 一方面是椭圆曲线及其模 $p$ 超奇异约化(当这一现象确实发生时), 另一方面是高度大于 $1$ 的 $p$-进整数环上的形式群.

我曾选修或旁听过 Tate 关于椭圆曲线算术的第一门课, 因而对这一切早有准备. 此外, 我知道 Weierstrass 预备定理, 以及它为任何掌握它的人所带来的威力. 在为我的学位论文证明某个结果时, 我曾考虑考察形式群上的挠点, 而且我认为对我而言这很清楚: 它们构成自同态环上的一个模. 请特别注意, 把它们用作 Galois 群的表示模完全不是我的想法.

但 Tate 时时刻刻都在我肩后看着, 无疑看到了许多我未曾考虑的事情. 在我提交学位论文时, 我还没有构造出高度为 $h$、自同态环 $\mathfrak{o}$ 等于 $\Q_p$ 上 $h$ 次局部域 $k$ 的整数环的形式群. 我只对非分歧情形有构造, 而且用的是基于 Lazard 技巧、极其繁琐的逐次次数方法. 在我完成学位论文后不久, 我坐在从 Brunswick 开往 Boston 的巴士上, 发现不仅能在所有情形构造出具有这种极大自同态结构的形式群, 而且其中有一个可以取多项式形式 $\pi x+x^{q}$. Tate 告诉我, 当他看到这一点时, 一切都各就各位了(Everything Fell Into Place). 其结果就是我们论文中那条精彩而优美的引理, 对此我完全不敢邀功. 我的记忆一向不可靠, 但我记得论文其余部分相当迅速地凝聚成形. 别忘了 Tate 当时已是类域论各个方面的大师. 可是, 如果你形式群的自同态环是 $\mathfrak{o}$, 而该形式群的 Tate 模是此自同态环上的秩一模, 那么 $k$ 上 $k(F[p^{\infty}])$ 的 Galois 群与子群 $\mathfrak{o}^{*}\subset k^{*}$ 之间的同构, 难道还不足以让你想到互反映射吗?
\end{namedblock}

\begin{namedblock}{TATE 致 SERRE 的信, 1964 年 1 月 10 日; COLLECTED WORKS, L8}
[由 Lubin 的结果]我们得到一个同态 $G_{K}\to U_{K}=K$ 的单位群. 把它限制到惯性群 $G_{K^{\mathrm{un}}}$ 上, 我们得到一个同态 $G_{K^{\mathrm{un}}}\to U_{K}$; 由上文的唯一性注记, 它很可能是典范的, 而且几乎无法怀疑它实际上就是\emph{互反律同构}(或者其负元!)\ldots 奇迹似乎在于: 一旦放弃代数群而转向形式群, 复乘理论便\emph{普遍地}(局部地)适用.
\end{namedblock}

\phantomsection
\section*{A 附录: 无限 Galois 理论与逆向极限}\label{sec:I.A}%
\addcontentsline{toc}{section}{A 附录: 无限 Galois 理论与逆向极限}
\index{wuxiangaloislilun@无限 Galois 理论(infinite Galois theory)}\index{nixiangjixian@逆向极限(inverse limit)}

我们回顾本章所需的两个专题 (也见 FT 第 7 章).

% 附录中定理按 A.x 编号
\renewcommand{\thetheorem}{A.\arabic{theorem}}
\setcounter{theorem}{0}

\subsection*{Galois 理论用于无限扩张}

固定一个域 $K$.

\begin{definition}\label{def:A.1}
一个域 $\Omega\supset K$ (未必有限次) 称为在 $K$ 上\emph{Galois}\index{galuoiskuozhang@Galois 扩张(Galois extension)!wuxian@无限(infinite)}, 如果

(a) 它在 $K$ 上代数且可分, 即 $\Omega$ 中每个元素都是系数在 $K$ 中的多项式的单根;

(b) 它在 $K$ 上\emph{正规}\index{zhengguikuozhang@正规扩张(normal extension)}, 即每个系数在 $K$ 中、在 $\Omega$ 中有根的不可约多项式在 $\Omega[X]$ 中分裂.
\end{definition}

\begin{proposition}\label{prop:A.2}
一个域 $\Omega$ 在 $K$ 上 Galois, 当且仅当它是 $K$ 的有限 Galois 扩张的并.
\end{proposition}

\begin{proof}
设 $\Omega$ 在 $K$ 上 Galois. 对每个 $\alpha\in\Omega$, $\alpha$ 在 $K$ 上的极小多项式在 $\Omega$ 中的分裂域是 $K$ 的有限 Galois 扩张, 而 $\Omega$ 是这些域的并. 反之, 若 $\Omega$ 是 $K$ 的有限 Galois 扩张的并, 则它在 $K$ 上代数且可分. 此外, 若 $f(X)\in K[X]$ 在 $\Omega$ 中有根, 则它在 $\Omega$ 的某个有限 Galois 子扩张 $E$ 中有根, 因此在 $E[X]$ 中分裂.
\end{proof}

若 $\Omega$ 是 $K$ 的 Galois 扩张, 则 Galois 群 $\Gal(\Omega/K)$ 定义为固定 $K$ 中元素的 $\Omega$ 的自同构群, 赋予如下拓扑: 以集合
\[
\Gal(\Omega/E),\quad \Omega\supset E\supset K,\ [E:K]<\infty
\]
作为 $1$ 处的基本邻域系. 这意味着 $\Gal(\Omega/K)$ 中两个元素接近, 是指它们在某个 “大” 的域 $E$ ($\Omega\supset E\supset K$, $[E:K]<\infty$) 上一致.

\begin{proposition}\label{prop:A.3}
群 $\Gal(\Omega/K)$ 是紧的且 Hausdorff 的.
\end{proposition}

\begin{proof}
考虑映射
\[
\sigma\mapsto\sigma|_E\colon\Gal(\Omega/K)\to\prod\Gal(E/K),
\]
其中乘积取遍所有 $E\subset\Omega$ 且在 $K$ 上有限 Galois 的子域. 命题 \ref{prop:A.2} 表明该映射是单的. 当我们给每个群 $\Gal(E/K)$ 赋予离散拓扑, 给乘积赋予积拓扑时, 在 $\Gal(\Omega/K)$ 上诱导的拓扑就是上面的拓扑. 我断言像是闭的. 像等于满足如下条件的族 $(\sigma_E)_E$ 的集合: 当 $E'\supset E$ 时 $\sigma_{E'}|_E=\sigma_E$. 设 $(\sigma_E)$ 不在像中. 则存在一对域 $E_2\supset E_1$ 使得 $\sigma_{E_2}|_{E_1}\neq\sigma_{E_1}$. 令
\[
U=\prod_{E\neq E_1,E_2}\Gal(E/K)\times\{\sigma_{E_2}\}\times\{\sigma_{E_1}\}.
\]
这是 $(\sigma_E)_E$ 的一个开邻域, 且 $U\cap\mathrm{Im}(\Gal(\Omega/K))=\varnothing$.
\end{proof}

$\Gal(\Omega/K)$ 上的拓扑离散, 当且仅当 $\Omega$ 是 $K$ 的有限扩张.

\begin{theorem}\label{thm:A.4}
设 $\Omega$ 是 $K$ 的 (可能无限的) Galois 扩张, Galois 群为 $G$. 则在 $\Omega$ 的子域与 $G$ 的闭子群之间存在一一对应. 更精确地:

(a) 对 $\Omega$ 的子域 $E$, $H\stackrel{\mathrm{def}}{=}\Gal(\Omega/E)$ 是 $G$ 的闭子群, 且 $E=\Omega^H$.

(b) 若 $H$ 是 $G$ 的子群, 则 $\Gal(\Omega/\Omega^H)$ 是 $H$ 在 $\Gal(\Omega/K)$ 中的闭包.

此外, $G$ 的正规闭子群对应于 $K$ 的 Galois 扩张, 而 $G$ 的开子群对应于 $K$ 的有限扩张.
\end{theorem}

\begin{proof}
设 $E\subset\Omega$ 是 $K$ 的有限扩张. 由于每个 $K$-同态 $E\to\Omega$ 都可扩张到 $\Omega$, 映射 $\sigma\mapsto\sigma|_{\Omega}\colon\Gal(\Omega/K)\to\Hom_K(E,\Omega)$ 是满的, 因此诱导一个双射
\[
\Gal(\Omega/K)/\Gal(\Omega/E)\to\Hom_K(E,\Omega).
\]
这表明 $\Gal(\Omega/E)$ 在 $\Gal(\Omega/E)$ 中的指数为 $[E:K]$. 由于它开 (按定义), 它也闭 (其补集是有限个开陪集的并).

设 $E\subset\Omega$ 是 $K$ 的任意扩张. 则 $E=\bigcup E_i$, 其中 $E_i$ 跑遍 $E$ 中包含的 $K$ 的有限扩张. 相应地, $\Gal(\Omega/E)=\bigcap\Gal(\Omega/E_i)$, 因此闭. 此外, 若 $\alpha\in\Omega$ 不被 $\Gal(\Omega/E)$ 固定, 则它不被任何 $\Gal(\Omega/E_i)$ 固定, 从而不在 $E$ 中. 于是 $E=\Omega^{\Gal(\Omega/E)}$.

设 $H$ 是 $\Gal(\Omega/K)$ 的子群, 并令 $H'=\Gal(\Omega/\Omega^H)$. 由有限扩张的 Galois 理论, 对每个 $\Gal(\Omega/K)$ 的开子群 $U$ 有 $UH=UH'$; 而由拓扑群理论, $H$ 的闭包是 $\bigcap UH$ (交取遍 $\Gal(\Omega/K)$ 的开子群 $U$), 因此 $\bar{H}=\bigcap UH=\bigcap UH'=H'$.
\end{proof}

\begin{example}\label{ex:A.5}
(a) 给 $\Z$ 赋予以有限指数子群为基本邻域系的拓扑, 并令 $\hat{\Z}$ 为完备化. 则 $\hat{\Z}=\prod_{\ell\text{ 素数}}\Z_{\ell}$, 且对每个 $m$ 有 $\hat{\Z}/m\hat{\Z}=\Z/m\Z$. 令 $\F$ 为 $\F_q$ 的代数闭包. 存在典范同构
\[
\hat{\Z}\to\Gal(\F/\F_q)
\]
把 $1\in\hat{\Z}$ 送到自同构 $x\mapsto x^q$. $\F_q$ 的 $m$ 次扩张对应于 $\hat{\Z}$ 的子群 $m\hat{\Z}$. 设 $\sigma$ 是 $\F/\F_q$ 的满足 $\sigma(x)=x^q$ 的自同构. 对 $\alpha\in\hat{\Z}$, 定义 $\sigma^{\alpha}$ 为 $\Gal(\F/\F_q)$ 中如下元素: 对每个 $m$, 当 $a\in\Z$ 充分接近 $\alpha$ (依赖于 $m$) 时 $\sigma^{\alpha}|_{\F_{q^m}}=\sigma^a|_{\F_{q^m}}$. 上述同构把 $\alpha$ 送到 $\sigma^{\alpha}$. (对 $\hat{\Z}$ 与同构 $\hat{\Z}\to\Gal(\F/\F_q)$ 的详细描述, 见 \cite{Artin1951}, 9.2.)

(b) 令 $\Omega_p=\bigcup_r\Q[\zeta_{p^r}]$, 其中 $\zeta_{p^r}$ 是一个本原 $p^r$ 次单位根. 对 $u\in\Z_p^{\times}$, 写 $u=a_0+a_1p+a_2p^2+\cdots$, $0\le a_i\le p-1$, 并定义
\[
\zeta_{p^r}^u=\zeta_{p^r}^{a_0+a_1p+\cdots+a_s p^s},\quad\text{每个 }s\ge r-1.
\]
这定义了 $\Z_p^{\times}$ 在 $\Omega_p$ 上的一个作用, 实际上是一个拓扑群同构
\[
\Z_p^{\times}\to\Gal(\Omega_p/\Q).
\]

(c) 令 $\Omega=\bigcup\Q[\zeta_n]$, 其中 $\zeta_n$ 是一个本原 $n$ 次单位根. 则
\[
\Omega=\Omega_2\cdot\Omega_3\cdot\Omega_5\cdots,\quad \Omega_p\cap\Omega_{p'}=\Q,\ p\neq p'.
\]
正如有限个有限 Galois 扩张的情形, 这推出
\[
\Gal(\Omega/\Q)=\prod\Gal(\Omega_p/\Q)
\]
(闭子群拓扑积). 于是存在一个同构
\[
\hat{\Z}^{\times}\to\Gal(\Omega/\Q).
\]
它可以描述如下: 若 $\zeta$ 是一个 $n$ 次单位根且 $u\in\Z^{\times}$, 则
\[
\zeta^u=\zeta^m\quad\text{对任何满足 }m\equiv u\bmod n\text{ 的 }m\in\Z.
\]

(d) 令 $\Omega_p=\bigcup\Q[\zeta_{p^r}]$, 如 (b). 则
\[
\Z_p^{\times}\simeq\Delta_p\times C_p,
\]
其中对 $p\neq 2$ 有 $\Delta_p\simeq(\Z/(p))^{\times}\approx\Z/(p-1)$, 而 $\Delta_2=\Z/2\Z$, 且 $C_p\approx\Z_p$. 令 $\Omega_p'=\Omega_p^{\Delta_p}$. 则 $\Gal(\Omega_p'/\Q_p)\approx\Z_p$. 令
\[
\Omega'=\Omega_2'\cdot\Omega_3'\cdots
\]
(在 $\Q^{\mathrm{al}}$ 内的复合). 则 $\Gal(\Omega'/\Q)\approx\prod\Z_p\simeq\hat{\Z}$.

(e) 对 $\Q$ 的每个有限扩张 $K$, $K\cdot\Omega'$ 是 $K$ 的 Galois 扩张, 其 Galois 群同构于 $\Z$ 中一个有限指数的子群. 但每个这样的子群又同构于 $\Z$ (因为它又是一个 $\Z$ 的有限指数子群的完备化). 因此 $\Gal(K\cdot\Omega'/K)\approx\Z$. 若我们固定一个同构, 并令 $K_m$ 为对应于 $m\Z$ 的域, 则我们看到:

(a) $K_m$ 是次数为 $m$ 的循环扩张;

(b) $K_m$ 是分圆的, 即包含在 $K$ 添加单位根所得的扩张中.
\end{example}

\begin{remark}\label{rem:A.6}
一般, 即使 $K=\Q$, 也可能存在 $\Gal(\Omega/K)$ 中有限指数但不开的子群 (见 FT, 7.26).
\end{remark}

\subsection*{逆向极限}

一个偏序集 $(I,\preceq)$ 称为\emph{有向的}\index{youxiangji@有向集(directed set)}, 如果对任意 $\alpha,\beta\in I$, 存在 $\gamma\in I$ 使得 $\alpha,\beta\preceq\gamma$. 例如, 以整除关系排序的正整数集
\[
m\preceq n\iff m\mid n
\]
是有向的.

一个\emph{逆向系}\index{nixiangxi@逆向系(inverse system)} (或投射系) 的集合族是指一族由有向集 $I$ 索引的集合 $(S_{\alpha})_{\alpha\in I}$, 连同对每个对 $\alpha\preceq\beta$ 的一个映射 $\varphi_{\alpha,\beta}\colon S_{\beta}\to S_{\alpha}$, 使得

(a) 对所有 $\alpha\in I$, $\varphi_{\alpha,\alpha}$ 是恒等映射;

(b) 对所有 $I$ 中的 $\alpha\preceq\beta\preceq\gamma$, $\varphi_{\alpha,\beta}\circ\varphi_{\beta,\gamma}=\varphi_{\alpha,\gamma}$.

一个集合 $S$ 连同对每个 $\alpha\in I$ 的一个映射 $\varphi_{\alpha}\colon S\to S_{\alpha}$ 使得 $\varphi_{\alpha}=\varphi_{\alpha,\beta}\circ\varphi_{\beta}$, 若满足明显的万有性质, 则称为逆向系 $(S_{\alpha})$, $(\varphi_{\alpha,\beta})$ 的\emph{逆向极限}\index{nixiangjixian@逆向极限(inverse limit)} (或投射极限).

每个集合、群或环的逆向系都有逆向极限. 例如 $\hat{\Z}=\varprojlim\Z/m\Z$.

\begin{example}\label{ex:A.7}
离散赋值环 $R$ 的完备化 $\hat{R}$ 典范同构于 $\varprojlim R/\mathfrak{m}^n$, 其中 $\mathfrak{m}$ 是 $R$ 中的极大理想.
\end{example}

一个群 $G$ 的\emph{profinite 完备化}\index{profinitewanchenghua@profinite 完备化(profinite completion)}定义为 $\varprojlim G/H$, 其中 $H$ 跑遍 $G$ 中有限指数的正规子群.

\begin{proposition}\label{prop:A.8}
有限 Abel 群的正合列组成的逆向系的逆向极限仍正合.
\end{proposition}

\begin{remark}\label{rem:A.9}
Galois 群 $\Gal(\Omega/K)$ 是群 $\Gal(E/K)$ 的投射极限, 其中 $E$ 跑遍 $\Omega$ 中在 $K$ 上有限且 Galois 的子域. 一个拓扑群若是有限群的一个投射极限, 则称为\emph{profinite 群}\index{profinitequn@profinite 群(profinite group)}. 它们恰为紧的全不连通拓扑群. (一个拓扑空间全不连通, 是指它的连通分支都是单点集.) 见 \cite{Serre1964} 或 \cite{Shatz1972}.
\end{remark}

% 恢复定理计数器为 节.序号 格式, 供后续章节使用
\renewcommand{\thetheorem}{\thesection.\arabic{theorem}}
